\documentclass[11pt]{article}

\usepackage[utf8]{inputenc}
% microtype 의 폰트 확장은 scalable 폰트를 요구한다. 기본 CM 이 비트맵(Type3)으로
% 잡히는 배포판에서는 lmodern 없이는 "auto expansion is only possible with
% scalable fonts" 로 컴파일이 죽는다.
\usepackage{lmodern}
\usepackage[T1]{fontenc}
\usepackage[margin=1in]{geometry}
\usepackage{amsmath,amssymb}
\usepackage{graphicx}
\usepackage{longtable,booktabs}
\usepackage{microtype}
% svgnames 를 줘야 teal 이 정의된다. 없으면 \cite 하나마다
% "Undefined color `teal'" 이 뜬다 (본편 산출물에서 732회).
\usepackage[svgnames]{xcolor}
% hyperref goes last; it makes \cite clickable through to the bibliography.
\usepackage[colorlinks=true,linkcolor=blue,citecolor=teal,urlcolor=magenta]{hyperref}

% pandoc emits \tightlist but expects the preamble to define it.
\providecommand{\tightlist}{%
  \setlength{\itemsep}{0pt}\setlength{\parskip}{0pt}}

% If the compile times out on Overleaf, comment out \tableofcontents below first.
\title{Visual Adversarial Attacks and Defenses in the Physical World: A Comprehensive Survey}
\author{Generated by AutoSurvey}
\date{}

\begin{document}
\maketitle
\tableofcontents
\newpage

\section{Introduction to Visual Adversarial Attacks}

\subsection{Definition and Purpose of Visual Adversarial Attacks}

Visual adversarial attacks refer to the process of crafting and applying subtle, often imperceptible, perturbations to input images in order to deceive computer vision systems into making incorrect predictions or classifications \cite{ref1}. These attacks can be launched in both digital and physical domains, with the latter being more realistic and threatening as the perturbations are introduced to the input before it is captured and converted into a binary image inside the vision system \cite{ref1}. The purpose of visual adversarial attacks is to exploit the vulnerabilities of deep learning models, which are widely used in various computer vision applications such as image classification, object detection, and facial recognition.

At its core, the concept of visual adversarial attacks is based on the idea that small, carefully crafted perturbations can be added to an image to alter its classification or detection by a computer vision system \cite{ref2}. These perturbations can be so subtle that they are imperceptible to the human eye, yet they can cause a deep learning model to misclassify an image with high confidence \cite{ref3}. The potential impact of visual adversarial attacks on computer vision systems is significant, as they can compromise the security and reliability of these systems in various applications \cite{ref4}.

Visual adversarial attacks can be categorized into different types, including physical and digital attacks \cite{ref5}. Physical attacks involve adding perturbations to the physical environment, such as printing a sticker on a stop sign, while digital attacks involve adding perturbations to the digital image \cite{ref1}. Both types of attacks can be effective in deceiving computer vision systems, but physical attacks are more realistic and threatening as they can be launched in the real world \cite{ref5}. The definition and purpose of visual adversarial attacks are closely related to the concept of adversarial examples, which are inputs to a machine learning model that are designed to cause the model to make a mistake \cite{ref4}.

The potential consequences of visual adversarial attacks are far-reaching and can have significant impacts on various applications, including autonomous vehicles, surveillance systems, and facial recognition systems \cite{ref6}. For instance, in the case of autonomous vehicles, visual adversarial attacks can be used to cause a vehicle to misclassify a stop sign or a pedestrian, which can have serious consequences \cite{ref6}. Similarly, in the case of facial recognition systems, visual adversarial attacks can be used to cause a system to misclassify an individual's identity, which can have serious consequences in applications such as security and law enforcement \cite{ref6}.

In order to develop effective defense mechanisms against visual adversarial attacks, it is essential to understand the different types of attacks and their characteristics \cite{ref4}. This understanding can inform the development of more robust and secure deep learning models, as well as the creation of detection and mitigation strategies \cite{ref3}. The study of visual adversarial attacks has the potential to lead to new insights and innovations in the field of computer vision, and it is essential to consider the potential risks and consequences of these attacks in order to develop effective defense mechanisms \cite{ref6}.

As the field of computer vision continues to evolve, the development of visual adversarial attacks will likely become increasingly sophisticated and powerful \cite{ref4}. Therefore, it is crucial to develop effective defense mechanisms to mitigate these risks and ensure the security and reliability of computer vision systems in various applications \cite{ref6}. By understanding the concept and purpose of visual adversarial attacks, as well as their potential consequences, researchers and developers can work towards creating more robust and secure computer vision systems that can withstand these types of attacks \cite{ref5}. Ultimately, the study of visual adversarial attacks has the potential to lead to significant advancements in the field of computer vision, and it is essential to continue researching and developing effective defense mechanisms against these attacks \cite{ref3}.

\subsection{Types of Visual Adversarial Attacks}

Visual adversarial attacks can be broadly categorized into three types: physical, digital, and hybrid attacks, each with its unique characteristics, advantages, and limitations. As discussed earlier, the concept of visual adversarial attacks is based on the idea that small, carefully crafted perturbations can be added to an image to alter its classification or detection by a computer vision system \cite{ref2}. These perturbations can be so subtle that they are imperceptible to the human eye, yet they can cause a deep learning model to misclassify an image with high confidence \cite{ref3}.

Physical adversarial attacks involve modifying the physical environment to deceive a computer vision system, such as adding stickers or patches to objects \cite{ref7}, or using light projections to create adversarial patterns \cite{ref8}. These attacks can be launched in the real world, making them more realistic and threatening as they can be designed to blend in with the surrounding environment \cite{ref9}. However, physical attacks can be limited by the need to physically access the target object or environment.

In contrast, digital adversarial attacks involve modifying the digital input to a computer vision system, such as adding noise or perturbations to images \cite{ref1}, or generating adversarial examples using machine learning algorithms \cite{ref10}. Digital attacks are often easier to launch than physical attacks, as they can be done remotely and do not require physical access to the target system. This highlights the importance of understanding the different types of visual adversarial attacks, as each type requires distinct defense strategies.

Hybrid adversarial attacks combine elements of both physical and digital attacks, offering a combination of the advantages of physical and digital attacks \cite{ref11}. These attacks can involve modifying the physical environment and then using digital techniques to enhance the attack, making them more effective than either physical or digital attacks alone \cite{ref12}. One of the key advantages of physical adversarial attacks is their ability to be stealthy and robust \cite{ref13}, while digital adversarial attacks have the advantage of being easier to launch and more flexible than physical attacks \cite{ref14}.

Understanding the different types of visual adversarial attacks is crucial for developing effective defense strategies and improving the security of computer vision systems \cite{ref15}. As the field of computer vision continues to evolve, the development of visual adversarial attacks will likely become increasingly sophisticated and powerful \cite{ref4}. Therefore, it is essential to continue researching and developing effective defense mechanisms against these attacks, which will be discussed in the following section. By recognizing the unique characteristics, advantages, and limitations of each type of attack, researchers and developers can work towards creating more robust and secure computer vision systems that can withstand these types of attacks \cite{ref16}.

\subsection{Threats and Risks of Visual Adversarial Attacks}

The threats and risks posed by visual adversarial attacks to deep learning models and computer vision systems are multifaceted and far-reaching, building upon the various types of attacks discussed earlier, including physical, digital, and hybrid attacks. One of the primary concerns is the potential for misclassification, where an adversarial example is crafted to cause a model to incorrectly classify an input \cite{ref17}. This can have significant consequences in applications such as autonomous vehicles, where incorrect classification of a stop sign or traffic light could lead to accidents \cite{ref18}. Furthermore, visual adversarial attacks can also be used to manipulate the output of a model, allowing an attacker to influence the decision-making process of a system \cite{ref19}.

As visual adversarial attacks can be highly sophisticated and targeted, they can exploit vulnerabilities in the computer vision system that are not immediately apparent, highlighting the need for robust defenses against these attacks. Another risk associated with visual adversarial attacks is disruption, where an attacker intentionally crafts inputs to cause a model to produce incorrect or misleading outputs \cite{ref20}. This can be particularly problematic in applications such as surveillance systems, where incorrect or misleading outputs could lead to false alarms or missed detections \cite{ref21}. Additionally, visual adversarial attacks can also be used to compromise the security of computer vision systems, allowing attackers to bypass security measures or gain unauthorized access to sensitive information \cite{ref22}.

The vulnerability of deep learning models to visual adversarial attacks is a significant concern, as these models are increasingly being used in a wide range of applications, including computer vision, natural language processing, and robotics \cite{ref4}. The use of deep learning models in safety-critical systems, such as autonomous vehicles and medical diagnosis, highlights the need for robust defenses against visual adversarial attacks \cite{ref23}. Furthermore, the ability of visual adversarial attacks to compromise the security of computer vision systems has significant implications for the development of secure and reliable AI systems \cite{ref24}.

To mitigate the risks associated with visual adversarial attacks, it is essential to develop robust defenses against these attacks, which will be discussed in the following section \cite{ref16}. This can include techniques such as adversarial training, where a model is trained on a dataset that includes adversarial examples \cite{ref25}. Additionally, techniques such as input preprocessing and detection methods can be used to identify and mitigate visual adversarial attacks \cite{ref26}. By developing and implementing effective defenses, we can mitigate the risks associated with visual adversarial attacks and ensure the secure and reliable operation of computer vision systems \cite{ref16}. Furthermore, ongoing research in this area is essential to stay ahead of the evolving threats and to develop more effective defenses against visual adversarial attacks \cite{ref27}.

\subsection{Real-World Implications of Visual Adversarial Attacks}

Visual adversarial attacks have significant real-world implications, particularly in applications that rely heavily on computer vision systems. As discussed earlier, the threats and risks posed by visual adversarial attacks are multifaceted and far-reaching, and it is essential to develop robust defenses against these attacks. One of the most critical areas of concern is autonomous vehicles, which use computer vision to perceive their surroundings and make decisions \cite{ref28}. Researchers have demonstrated that adversarial attacks can be used to manipulate traffic signs, causing autonomous vehicles to misinterpret them and potentially leading to accidents \cite{ref29}. For example, an attacker could create a physical adversarial example that causes an autonomous vehicle to misrecognize a stop sign as a speed limit sign, leading to a potential collision \cite{ref30}.

In addition to autonomous vehicles, visual adversarial attacks also have significant implications for surveillance systems, which rely on computer vision to detect and track individuals \cite{ref11}. Adversarial attacks can be used to evade detection or manipulate the system's output, allowing attackers to avoid detection or gain unauthorized access to secure areas. Facial recognition systems are also vulnerable to visual adversarial attacks, which can be used to manipulate the system's output and potentially lead to misidentification \cite{ref31}. This has significant implications for applications such as border control, law enforcement, and access control, where facial recognition systems are used to verify identities. For instance, an attacker could create a physical adversarial example that causes a facial recognition system to misrecognize them as someone else, allowing them to gain unauthorized access to secure areas \cite{ref32}.

The real-world implications of visual adversarial attacks are not limited to these examples. Any application that relies on computer vision is potentially vulnerable to these types of attacks, and the consequences can be severe. Therefore, it is essential to develop effective defenses against visual adversarial attacks to ensure the security and reliability of computer vision systems \cite{ref4}. To address this challenge, researchers have proposed various defense mechanisms against visual adversarial attacks, including adversarial training, input preprocessing, and detection methods \cite{ref33}. These defense mechanisms will be discussed in more detail in the following section, which will explore the current state of defenses against visual adversarial attacks and propose future directions for research.

In the context of developing effective defenses, it is essential to consider the sophisticated and targeted nature of visual adversarial attacks. Adversarial attacks can exploit vulnerabilities in the computer vision system that are not immediately apparent, making it challenging to develop robust defenses. Therefore, it is crucial to continue researching and developing new defense mechanisms against visual adversarial attacks to stay ahead of potential attackers \cite{ref16}. By developing and implementing effective defenses, we can mitigate the risks associated with visual adversarial attacks and ensure the secure and reliable operation of computer vision systems.

\section{Background and Fundamentals of Adversarial Attacks}

\subsection{Introduction to Adversarial Attacks}

Adversarial attacks have emerged as a significant threat to the security and reliability of computer vision systems, which are increasingly being used in various applications such as self-driving cars, surveillance, and facial recognition. The concept of adversarial attacks was first introduced in the context of image classification, where it was shown that adding small amounts of noise to an image could cause a deep neural network to misclassify it \cite{ref4}. Since then, adversarial attacks have been extended to other computer vision tasks such as object detection, segmentation, and generation. These attacks involve crafting small, carefully designed perturbations to the input images that can deceive deep learning models into making incorrect predictions.

There are several types of adversarial attacks, including white-box, black-box, and gray-box attacks, which differ in the level of access the attacker has to the model's architecture, weights, and training data. White-box attacks assume that the attacker has full access to the model's architecture, weights, and training data, while black-box attacks assume that the attacker only has access to the input and output of the model \cite{ref34}. Gray-box attacks fall somewhere in between, where the attacker has some knowledge of the model's architecture or weights. Additionally, adversarial attacks can be categorized into targeted and untargeted attacks, where targeted attacks aim to cause the model to output a specific incorrect class or label, and untargeted attacks aim to cause the model to output any incorrect class or label \cite{ref35}.

The threats posed by adversarial attacks are significant, as they can be used to compromise the security and reliability of computer vision systems, with potentially catastrophic consequences. For example, an attacker could use an adversarial attack to cause a self-driving car to misrecognize a stop sign or a pedestrian \cite{ref5}. Similarly, an attacker could use an adversarial attack to compromise the security of a facial recognition system, allowing them to gain unauthorized access to secure facilities or systems \cite{ref36}. The threats posed by adversarial attacks are not limited to computer vision systems, as they can also be used to compromise the security of other machine learning systems, such as natural language processing systems \cite{ref37}.

To defend against adversarial attacks, several defense mechanisms have been proposed, including adversarial training, input preprocessing, and robust optimization \cite{ref6}. Adversarial training involves training the model on adversarial examples, which can help to improve the model's robustness to these attacks. Input preprocessing involves modifying the input data to reduce the effectiveness of adversarial attacks, such as by adding noise or using image compression. Robust optimization involves optimizing the model's parameters to minimize the loss function under adversarial attacks. Despite these defense mechanisms, adversarial attacks remain a significant threat to the security and reliability of computer vision systems.

Therefore, it is essential to continue researching and developing new defense mechanisms to protect against these attacks. One potential approach is to use techniques such as attention mechanisms and log-polar mapping to improve the robustness of computer vision models \cite{ref38}. Another approach is to use techniques such as generative models and autoencoders to detect and defend against adversarial attacks \cite{ref26}. By understanding the types of adversarial attacks and developing effective defense mechanisms, we can improve the security and reliability of computer vision systems and mitigate the threats posed by these attacks \cite{ref39}.

\subsection{Threat Models and Attack Vectors}

Threat models and attack vectors are crucial components in the study of adversarial attacks, as they define the capabilities and goals of the adversary. In the context of adversarial attacks, threat models can be categorized into three primary types: white-box, black-box, and gray-box attacks. Each of these threat models presents a different level of access to the target model, ranging from full access to no access at all. Understanding these threat models is essential to evaluating the robustness of machine learning models and developing effective defense mechanisms against adversarial attacks.

White-box attacks assume that the adversary has complete knowledge of the target model, including its architecture, parameters, and training data \cite{ref40}. This level of access allows the adversary to craft highly effective adversarial examples, often using gradient-based methods to optimize the perturbation \cite{ref41}. White-box attacks are typically used as a benchmark to evaluate the robustness of machine learning models, as they represent the most powerful type of attack. However, they are less realistic in practice, as an attacker is unlikely to have complete access to the target model. Despite this, white-box attacks provide a useful upper bound on the vulnerability of a model to adversarial attacks.

In contrast, black-box attacks assume that the adversary has no knowledge of the target model, including its architecture, parameters, and training data \cite{ref42}. In this scenario, the adversary must rely on indirect methods to craft adversarial examples, such as querying the target model and observing its outputs \cite{ref43}. Black-box attacks are more realistic in practice, as they reflect the typical scenario where an attacker has limited access to the target model. However, they are often less effective than white-box attacks, as the adversary must rely on indirect methods to craft adversarial examples. Black-box attacks are particularly relevant in scenarios where the model is deployed in a real-world setting, and the attacker must interact with the model through a limited interface.

Gray-box attacks fall between white-box and black-box attacks, assuming that the adversary has partial knowledge of the target model \cite{ref44}. This partial knowledge can include information about the model's architecture, parameters, or training data, but not all of it. Gray-box attacks are often used to evaluate the robustness of machine learning models in scenarios where the adversary has some knowledge of the target model, but not complete access. Gray-box attacks can provide a more realistic estimate of the vulnerability of a model to adversarial attacks, as they reflect the scenario where an attacker has some knowledge of the model, but not complete access.

In addition to these threat models, attack vectors are also an important consideration in adversarial attacks. Attack vectors refer to the specific methods used by the adversary to craft and deliver adversarial examples \cite{ref45}. Common attack vectors include gradient-based methods, optimization-based methods, and generative model-based methods \cite{ref46}. Each of these attack vectors has its own strengths and weaknesses, and the choice of attack vector often depends on the specific threat model and the goals of the adversary. The choice of attack vector can significantly impact the effectiveness of adversarial attacks, and understanding the different attack vectors is essential to developing effective defense mechanisms.

The choice of threat model and attack vector can significantly impact the effectiveness of adversarial attacks. For example, white-box attacks using gradient-based methods can be highly effective, but require complete access to the target model \cite{ref47}. In contrast, black-box attacks using query-based methods can be less effective, but require only limited access to the target model \cite{ref48}. Gray-box attacks, on the other hand, can offer a balance between effectiveness and accessibility, as they can be used to evaluate the robustness of machine learning models in scenarios where the adversary has partial knowledge of the target model \cite{ref49}. By considering the different threat models and attack vectors, researchers can develop more realistic and effective attacks, and developers can design more robust and secure machine learning models.

In recent years, there has been a growing interest in developing more sophisticated threat models and attack vectors, such as transfer-based attacks and ensemble-based attacks \cite{ref50}. These attacks can be used to evaluate the robustness of machine learning models in more realistic scenarios, where the adversary has limited access to the target model and must rely on indirect methods to craft adversarial examples. Additionally, there has been a growing interest in developing more effective defense mechanisms against adversarial attacks, such as adversarial training and input preprocessing \cite{ref51}. By understanding the different threat models and attack vectors, researchers can develop more effective defense mechanisms, and developers can design more robust and secure machine learning models \cite{ref52}.

\subsection{Common Techniques for Generating Adversarial Examples}

Generating adversarial examples is a crucial step in evaluating the robustness of machine learning models, as it allows developers and researchers to test the vulnerability of their models to different types of attacks. As discussed in the previous section, threat models and attack vectors play a significant role in the development of adversarial attacks, and the choice of technique used to generate adversarial examples depends on the specific use case and the characteristics of the model being attacked. In this subsection, we will discuss the common techniques used to generate adversarial examples, including gradient-based methods, optimization-based methods, and generative model-based methods.

Gradient-based methods are one of the most popular techniques used to generate adversarial examples. These methods utilize the gradient of the loss function to find the direction in which the input should be perturbed to maximize the loss. One of the earliest and most well-known gradient-based methods is the Fast Gradient Sign Method (FGSM) \cite{ref53}. FGSM generates adversarial examples by adding a perturbation to the input in the direction of the gradient of the loss function. Another popular gradient-based method is the Projected Gradient Descent (PGD) method \cite{ref54}, which generates adversarial examples by iteratively adding a perturbation to the input in the direction of the gradient of the loss function and then projecting the result back onto the feasible set.

In addition to gradient-based methods, optimization-based methods are another popular technique used to generate adversarial examples. These methods formulate the generation of adversarial examples as an optimization problem, where the goal is to find the minimum perturbation that causes the model to misclassify the input. One of the most well-known optimization-based methods is the Carlini and Wagner (C\&W) method \cite{ref55}, which generates adversarial examples by solving a box-constrained optimization problem. Another popular optimization-based method is the Momentum Iterative Fast Gradient Sign Method (MI-FGSM) \cite{ref56}, which generates adversarial examples by iteratively adding a perturbation to the input in the direction of the gradient of the loss function and then using momentum to stabilize the update.

More recently, generative model-based methods have been developed to generate adversarial examples. These methods utilize generative models, such as Generative Adversarial Networks (GANs) or Variational Autoencoders (VAEs), to generate adversarial examples. One of the most well-known generative model-based methods is the Generative Adversarial Perturbations (GAP) method \cite{ref57}, which generates adversarial examples by using a GAN to generate perturbations that cause the model to misclassify the input. Another popular generative model-based method is the Adversarial Example Generation (AEG) method \cite{ref58}, which generates adversarial examples by using a VAE to generate perturbations that cause the model to misclassify the input.

The choice of technique used to generate adversarial examples depends on the specific use case and the characteristics of the model being attacked. For example, gradient-based methods are often faster and more efficient than optimization-based methods, but may not be as effective at generating adversarial examples that are transferable across different models. On the other hand, optimization-based methods can generate more transferable adversarial examples, but may be slower and more computationally expensive. Generative model-based methods can generate highly realistic adversarial examples, but may require a large amount of training data and computational resources.

Furthermore, recent studies have also explored the use of evolutionary algorithms \cite{ref59} and particle swarm optimization \cite{ref60} to generate adversarial examples. These methods have shown promising results in generating transferable adversarial examples, but require careful tuning of hyperparameters and can be computationally expensive. Additionally, the use of adversarial training \cite{ref61} has been shown to be effective in improving the robustness of machine learning models. Adversarial training involves training the model on a mixture of clean and adversarial examples, which can help the model learn to be more robust to attacks.

In conclusion, generating adversarial examples is a critical step in evaluating the robustness of machine learning models. Various techniques have been developed to generate these examples, each with its strengths and weaknesses. By understanding these techniques and their limitations, researchers and practitioners can develop more effective methods for generating adversarial examples and improving the robustness of machine learning models. The following section will discuss the defenses against adversarial attacks, including the use of adversarial training, input preprocessing, and other techniques to improve the robustness of machine learning models \cite{ref62}.

\section{Taxonomy of Visual Adversarial Attacks}

\subsection{Physical Adversarial Attacks}

Physical adversarial attacks are a type of attack that involves manipulating the physical environment to deceive computer vision systems, which is a natural extension of the digital adversarial attacks discussed earlier. These attacks are particularly concerning because they can be used to compromise the security of systems that rely on computer vision, such as self-driving cars, surveillance systems, and facial recognition systems. In this subsection, we will introduce the concept of physical adversarial attacks, their types, and the threats they pose to computer vision systems, highlighting the key differences and similarities with digital adversarial attacks.

Physical adversarial attacks can be categorized into two main types: invasive and non-invasive attacks \cite{ref63}. Invasive attacks involve physically modifying the object or environment being captured by the camera, such as adding a sticker to a stop sign, which can be seen as a physical equivalent of adding noise or perturbations to a digital image. Non-invasive attacks, on the other hand, involve manipulating the light or other environmental factors to deceive the camera, such as using a projector to project an adversarial image onto a surface, which can be viewed as a physical manifestation of digital adversarial attacks.

One of the most significant threats posed by physical adversarial attacks is their ability to compromise the security of computer vision systems in the physical world, which can have serious consequences in applications such as self-driving cars, surveillance systems, and facial recognition systems \cite{ref5}. For example, an attacker could use a physical adversarial attack to manipulate a self-driving car's perception system, causing it to misinterpret its surroundings and potentially leading to an accident. Similarly, an attacker could use a physical adversarial attack to evade detection by a surveillance system, potentially allowing them to commit a crime without being detected \cite{ref11}.

Physical adversarial attacks can also be used to compromise the security of facial recognition systems, which is a critical application of computer vision \cite{ref32}. For instance, an attacker could use a physical adversarial attack to manipulate a facial recognition system's perception of their face, potentially allowing them to gain access to secure areas or commit identity theft. This type of attack is particularly concerning because facial recognition systems are increasingly being used in security-critical applications, such as border control and law enforcement.

In addition to their potential to compromise the security of computer vision systems, physical adversarial attacks also pose a significant challenge to the development of robust computer vision models \cite{ref13}. Because physical adversarial attacks can be used to manipulate the physical environment in a way that is not easily detectable by humans, they can be used to create adversarial examples that are highly effective at deceiving computer vision models. This can make it difficult to develop computer vision models that are robust to physical adversarial attacks, because the models must be able to detect and respond to a wide range of potential attacks.

To address the threat posed by physical adversarial attacks, researchers have proposed a number of different defense strategies, including adversarial training, input preprocessing, and object detection and tracking \cite{ref39}, \cite{ref64}, and \cite{ref65}. These defense strategies can be seen as an extension of the defense strategies used to defend against digital adversarial attacks, and they highlight the need for a comprehensive approach to defending against adversarial attacks in both the digital and physical domains.

In conclusion, physical adversarial attacks pose a significant threat to the security of computer vision systems in the physical world, and they highlight the need for a comprehensive approach to defending against adversarial attacks in both the digital and physical domains \cite{ref4}. By understanding the types of physical adversarial attacks that are possible, and the techniques that can be used to defend against them, researchers and practitioners can help to ensure the security and reliability of computer vision systems in a wide range of applications.

\subsection{Digital Adversarial Attacks}

Digital adversarial attacks are a type of cyber threat that involves manipulating digital images to deceive computer vision systems, posing a significant risk to the security and reliability of these systems in various applications, including self-driving cars, surveillance systems, and facial recognition systems \cite{ref66}. As a fundamental concept in the realm of adversarial attacks, digital adversarial attacks serve as a precursor to the more complex and sophisticated attacks that will be discussed in subsequent sections, such as physical and hybrid adversarial attacks \cite{ref6}.

One of the key challenges in defending against digital adversarial attacks is that they can be highly sophisticated and difficult to detect, often requiring significant computational power and complex algorithms to generate \cite{ref2}. For instance, an attacker may use a technique called ``gradient-based optimization'' to create an adversarial image that is specifically designed to maximize the error rate of the computer vision system. This can involve using large amounts of computational power to generate the adversarial image, making it difficult for defenders to keep up. Furthermore, digital adversarial attacks can be highly transferable, meaning that an attack that is designed to work against one computer vision system can also be effective against other systems \cite{ref67}.

To defend against digital adversarial attacks, several techniques can be employed, including ``adversarial training,'' which involves training the computer vision system on a dataset that includes adversarial images \cite{ref68}. This can help the system to learn to recognize and resist adversarial attacks, making it more robust and secure. Another approach is to use ``input preprocessing'' techniques, such as filtering or compression, to remove or reduce the effects of adversarial perturbations \cite{ref69}. Additionally, ``diversity'' techniques, such as using multiple computer vision systems or algorithms to analyze the same image, can make it more difficult for an attacker to launch a successful attack \cite{ref37}.

Recent studies have highlighted the potential risks and vulnerabilities of computer vision systems to digital adversarial attacks, including \cite{ref36}, which proposed a method for generating invisible adversarial examples that can fool face recognition systems, and \cite{ref70}, which proposed a method for generating imperceptible and transferable adversarial examples that can fool face recognition systems. These studies demonstrate the need for defenders to be aware of the risks posed by digital adversarial attacks and to take steps to mitigate them. As the field of computer vision continues to evolve, it is essential to develop effective defense mechanisms against digital adversarial attacks to ensure the security and reliability of computer vision systems in a wide range of applications \cite{ref71}.

In conclusion, digital adversarial attacks are a serious threat to computer vision systems, and defenders must be prepared to use a range of technical and strategic techniques to defend against them. By understanding the types of attacks that are possible, and the techniques that can be used to defend against them, defenders can help to ensure the security and reliability of computer vision systems in various applications, ultimately paving the way for the development of more robust and secure systems that can withstand the increasingly sophisticated threats posed by physical and hybrid adversarial attacks \cite{ref72}.

\subsection{Hybrid Adversarial Attacks}

Hybrid adversarial attacks combine the characteristics of physical and digital attacks, posing a significant threat to computer vision systems. These attacks leverage the strengths of both physical and digital attacks, making them more potent and challenging to defend against. In this subsection, we will delve into the concept of hybrid adversarial attacks, their types, and the threats they pose to computer vision systems, building on the foundation established in the previous discussion on digital adversarial attacks.

As we previously discussed, digital adversarial attacks can be highly sophisticated and difficult to detect, and the emergence of hybrid adversarial attacks further complicates the defense landscape. Hybrid adversarial attacks can be categorized into two main types: physical-digital hybrid attacks and digital-physical hybrid attacks. Physical-digital hybrid attacks involve creating physical adversarial examples that are then digitized and processed by computer vision systems \cite{ref63}. For instance, an attacker can create a physical adversarial patch that is then captured by a camera and processed by a computer vision system. On the other hand, digital-physical hybrid attacks involve creating digital adversarial examples that are then printed or displayed in the physical world \cite{ref32}.

The threats posed by hybrid adversarial attacks are significant, as they can exploit the vulnerabilities of both physical and digital systems. For example, an attacker can create a physical adversarial example that is designed to evade detection by digital systems, but still causes misclassification by the computer vision system \cite{ref73}. Additionally, hybrid adversarial attacks can be used to launch targeted attacks, where the attacker aims to misclassify a specific object or scene \cite{ref65}. This highlights the need for defense mechanisms that can effectively counter both physical and digital components of hybrid adversarial attacks.

The emergence of hybrid adversarial attacks has significant implications for the security and reliability of computer vision systems. As computer vision systems become increasingly ubiquitous in various applications, including autonomous vehicles, surveillance systems, and facial recognition systems, the risk of hybrid adversarial attacks also increases \cite{ref4}. Therefore, it is essential to develop effective defense mechanisms against hybrid adversarial attacks to ensure the security and reliability of computer vision systems. Several defense mechanisms have been proposed to counter hybrid adversarial attacks, including adversarial training, input preprocessing, and robust optimization \cite{ref34}. Adversarial training involves training computer vision systems on adversarial examples to improve their robustness against attacks. Input preprocessing involves modifying the input data to reduce the effectiveness of adversarial attacks. Robust optimization involves optimizing the computer vision system's parameters to improve its robustness against attacks.

However, developing effective defense mechanisms against hybrid adversarial attacks is a challenging task. Hybrid adversarial attacks can be highly sophisticated and adaptive, making it difficult to develop defense mechanisms that can detect and mitigate them \cite{ref7}. Additionally, the complexity of computer vision systems and the variability of hybrid adversarial attacks make it challenging to develop defense mechanisms that can generalize across different scenarios and applications. As we will discuss in the following section, the development of effective defense mechanisms against hybrid adversarial attacks requires a comprehensive understanding of the attack landscape and the development of novel defense strategies.

In conclusion, hybrid adversarial attacks pose a significant threat to computer vision systems, and developing effective defense mechanisms against them is essential to ensure the security and reliability of these systems. By understanding the characteristics and types of hybrid adversarial attacks, researchers and practitioners can develop more effective defense mechanisms to counter these attacks. Further research is needed to develop more sophisticated defense mechanisms that can detect and mitigate hybrid adversarial attacks in various applications and scenarios \cite{ref37}.

\section{Methods for Generating Physical Adversarial Examples}

\subsection{Introduction to Physical Adversarial Examples}

Physical adversarial examples are a type of attack that can be used to deceive machine learning models in the physical world, posing a significant threat to the security and reliability of these systems. Unlike digital adversarial examples, which are designed to be effective in a digital environment, physical adversarial examples are designed to be effective in the real world, where they can be used to attack systems such as self-driving cars, surveillance cameras, and facial recognition systems \cite{ref74}. These attacks can be used to cause misclassification, manipulation, or disruption of the system, and can have serious consequences in safety-critical applications.

There are several types of physical adversarial examples, including physical adversarial patches, which are stickers or other objects that can be placed on a surface to deceive a machine learning model \cite{ref75}. Physical adversarial objects are another type, which are designed to be misclassified by a machine learning model \cite{ref76}. These physical adversarial examples can be designed to be highly effective, even in the presence of noise or other environmental factors \cite{ref77}.

Generating effective physical adversarial examples is a challenging task, as they must be designed to be effective in the real world, where they will be subject to various environmental factors such as lighting, noise, and weather \cite{ref63}. To address this challenge, researchers have developed various methods for generating physical adversarial examples, including optimization-based methods and generative models \cite{ref78}. These methods can be used to generate physical adversarial examples that are highly effective, even in the presence of noise or other environmental factors.

Ensuring the robustness of physical adversarial examples to various environmental factors is a key challenge in generating them. To address this challenge, researchers have developed methods for generating physical adversarial examples that are designed to be robust to noise, weather, and other environmental factors \cite{ref79}. Additionally, ensuring that physical adversarial examples are visually natural and do not raise suspicion is another challenge. To address this challenge, researchers have developed methods for generating physical adversarial examples that are designed to be visually natural and do not raise suspicion \cite{ref80}.

Evaluating the effectiveness of physical adversarial examples is also an important aspect of generating them. To address this challenge, researchers have developed methods for evaluating the effectiveness of physical adversarial examples, including metrics such as success rate, perturbation size, and robustness to noise \cite{ref81}. These metrics can be used to evaluate the effectiveness of physical adversarial examples and to compare the performance of different methods for generating them.

The development of physical adversarial examples has significant implications for the security and reliability of machine learning systems in the physical world. As machine learning systems become increasingly ubiquitous in safety-critical applications, the risk of physical adversarial attacks will become more significant \cite{ref5}. Therefore, it is essential to develop effective methods for generating and evaluating physical adversarial examples, as well as for defending against them. The study of physical adversarial examples is an active area of research, and there are many open questions and challenges that remain to be addressed \cite{ref7}.

In the next section, we will discuss the methods for generating physical adversarial patches, which are a crucial step in creating effective attacks against computer vision systems in the physical world \cite{ref82}.

\subsection{Methods for Generating Physical Adversarial Patches}

Generating physical adversarial patches is a crucial step in creating effective attacks against computer vision systems in the physical world. As discussed in the previous section, physical adversarial examples can be used to deceive machine learning models in various applications, including self-driving cars, surveillance cameras, and facial recognition systems \cite{ref74}. To generate such patches, various methods have been proposed, including optimization-based methods and generative models. Optimization-based methods, such as \cite{ref83}, utilize a two-step optimization architecture to optimize the location and size of adversarial patches. These methods have been shown to be effective in generating stealthy and efficient patches.

In addition to optimization-based methods, generative models have also been widely used for generating physical adversarial patches. For example, \cite{ref84} proposes a Dual-Perception-Based Framework (DPBF) to generate more vivid patches that can enhance transferability, stealthiness, and practicality. Another approach is to use a generative adversarial network (GAN) to generate physical adversarial patches. For instance, \cite{ref85} proposes a physically realizable and transferable adversarial patch framework tailored for multimodal large language models-based autonomous driving systems.

Other techniques have also been proposed to improve the transferability and effectiveness of adversarial patches. For example, \cite{ref86} has been proposed to improve the transferability of adversarial patches. The use of reinforcement learning has also been explored for generating physical adversarial patches. For example, \cite{ref87} proposes a method that utilizes a reinforcement learning framework to simultaneously optimize the position and perturbation of an adversarial patch. Furthermore, \cite{ref88} proposes an approach to generate adversarial yet inconspicuous patches with a single image. Other methods, such as \cite{ref89}, have been proposed to generate sparse and patch adversarial attacks.

The evaluation of physical adversarial patches is also an important aspect of generating effective attacks. Methods such as \cite{ref90} have been proposed to evaluate the effectiveness of physical adversarial patches. In addition to the methods mentioned above, other techniques such as \cite{ref91} have been proposed to generate physical adversarial patches. These methods will be further discussed in the next section, which will focus on evaluating and comparing different physical adversarial attack methods, including their success rate, perturbation size, and robustness to environmental conditions. Overall, the generation of physical adversarial patches is a complex task that requires careful consideration of various factors, including the desired properties of the patches, the specific use case, and the evaluation metrics.

\subsection{Evaluation and Comparison of Physical Adversarial Attack Methods}

Evaluating and comparing different physical adversarial attack methods is crucial to understanding their effectiveness and limitations. Various metrics can be used to assess the performance of these methods, including success rate, perturbation size, and robustness to different environmental conditions. The success rate of an attack method refers to the percentage of times it is able to successfully deceive the target model \cite{ref5}. Perturbation size, on the other hand, measures the amount of noise or distortion added to the input data to create the adversarial example \cite{ref92}. Robustness to environmental conditions, such as lighting, weather, or viewpoint changes, is also an essential aspect of physical adversarial attacks \cite{ref8}.

To facilitate a comprehensive comparison of physical adversarial attack methods, several studies have proposed benchmarking frameworks. For example, \cite{ref93} proposed a benchmarking framework for evaluating physical attacks against object detection models. The framework includes various metrics, such as success rate, average precision, and robustness to different environmental conditions. The study found that the success rate of physical attacks can vary significantly depending on the attack method, object detection model, and environmental conditions. Similarly, \cite{ref94} evaluated the robustness of vehicle detection models to physical adversarial attacks, using a dataset of images captured in various environmental conditions.

In addition to success rate and robustness, perturbation size is also an important metric for evaluating physical adversarial attack methods. A smaller perturbation size generally indicates a more effective attack method, as it requires less noise or distortion to deceive the target model \cite{ref95}. However, reducing perturbation size can also make the attack method more sensitive to environmental conditions, such as lighting or viewpoint changes. To address these challenges, researchers have proposed various methods for generating physical adversarial examples with small perturbation sizes, such as \cite{ref96}, which uses a gradient-based optimization algorithm to find the smallest perturbation that can deceive the target model.

The development of defense methodologies that can improve the robustness of machine learning models to physical adversarial attacks is also crucial for ensuring their security and reliability in real-world applications. For example, \cite{ref97} proposed a defense methodology that uses a combination of data augmentation techniques and adversarial training to improve the robustness of CNN models to physical adversarial attacks. The study found that the proposed methodology can significantly improve the robustness of CNN models to physical adversarial attacks, while also reducing the computational cost. Furthermore, \cite{ref98} highlights the importance of developing robust machine learning models that can withstand physical adversarial attacks.

In the context of camera-based smart systems, physical adversarial attacks can be used to deceive object detection models, which can have significant consequences in safety-critical applications, such as autonomous vehicles or surveillance systems \cite{ref99}. Therefore, it is essential to develop defense methodologies that can improve the robustness of these models to physical adversarial attacks. One approach is to use data augmentation techniques, such as random cropping and flipping, to improve the robustness of object detection models to environmental conditions. Another approach is to use adversarial training, which involves training the model on a dataset of adversarial examples \cite{ref100}. This can help improve the robustness of the model to physical adversarial attacks, but it requires a large dataset of adversarial examples, which can be challenging to generate.

In conclusion, evaluating and comparing different physical adversarial attack methods is essential for understanding their effectiveness and limitations. By using a comprehensive set of metrics, including success rate, perturbation size, and robustness to environmental conditions, researchers can develop more effective physical adversarial attack methods and improve the security and reliability of machine learning models in real-world applications. The development of defense methodologies that can improve the robustness of these models to physical adversarial attacks is crucial for ensuring their security and reliability in safety-critical applications.

\section{Defense Mechanisms Against Visual Adversarial Attacks}

\subsection{Introduction to Defense Mechanisms}

The concept of defense mechanisms against visual adversarial attacks is of paramount importance in the field of computer vision and machine learning. Visual adversarial attacks pose a significant threat to the security and reliability of deep learning models, which are widely used in various applications such as image classification, object detection, and facial recognition. The importance of defense mechanisms against visual adversarial attacks cannot be overstated, as they are crucial in preventing the misuse of deep learning models and ensuring the safety and security of sensitive information \cite{ref101}.

One of the primary challenges in developing effective defense mechanisms against visual adversarial attacks is the complexity and diversity of these attacks. Visual adversarial attacks can take many forms, including pixel-level attacks, object-level attacks, and scene-level attacks \cite{ref102}. Each type of attack requires a different defense strategy, and developing a single defense mechanism that can effectively counter all types of visual adversarial attacks is a significant challenge. Furthermore, the rapid evolution of visual adversarial attacks makes it essential to continuously update and improve defense mechanisms to stay ahead of potential threats \cite{ref103}.

The goals of defense mechanisms against visual adversarial attacks are multifaceted. Firstly, they aim to detect and prevent visual adversarial attacks from compromising the security and reliability of deep learning models. Secondly, they seek to improve the robustness of deep learning models against visual adversarial attacks, ensuring that they can maintain their performance and accuracy even in the presence of attacks \cite{ref104}. Finally, defense mechanisms against visual adversarial attacks aim to provide a framework for evaluating and comparing the effectiveness of different defense strategies, enabling researchers and practitioners to develop more effective and efficient defense mechanisms against visual adversarial attacks \cite{ref105}.

To address these challenges, several defense mechanisms against visual adversarial attacks have been proposed in recent years, each with its strengths and weaknesses. Adversarial training is a popular defense mechanism that involves training deep learning models on adversarial examples to improve their robustness against attacks \cite{ref106}. Input preprocessing is another defense mechanism that involves preprocessing the input data to remove or reduce the effectiveness of visual adversarial attacks \cite{ref107}. Other defense mechanisms include gradient masking, certified defenses, and robust optimization \cite{ref6}.

In addition to these techniques, novel approaches such as reinforcement learning and generative models have shown significant promise in improving the robustness of deep learning models against visual adversarial attacks \cite{ref108}. The use of these techniques can provide a framework for evaluating and comparing the effectiveness of different defense mechanisms, enabling researchers and practitioners to develop more effective and efficient defense mechanisms against visual adversarial attacks. Furthermore, the development of more effective and efficient defense mechanisms against visual adversarial attacks will require a deeper understanding of the types of attacks and the strengths and weaknesses of different defense strategies.

In conclusion, defense mechanisms against visual adversarial attacks are essential in preventing the misuse of deep learning models and ensuring the safety and security of sensitive information. The development of effective defense mechanisms against visual adversarial attacks is a complex and challenging task, requiring a deep understanding of the types of attacks and the strengths and weaknesses of different defense strategies. By providing a framework for evaluating and comparing the effectiveness of different defense mechanisms, researchers and practitioners can develop more effective and efficient defense mechanisms against visual adversarial attacks \cite{ref109}. This will ultimately lead to the development of robust and effective defense mechanisms against visual adversarial attacks, which is critical in ensuring the reliability and security of deep learning models in various applications. The next section will delve into the specifics of adversarial training and robustness certification, two crucial defense mechanisms against visual adversarial attacks.

\subsection{Adversarial Training and Robustness Certification}

Adversarial training and robustness certification are two crucial defense mechanisms against visual adversarial attacks, building on the foundation of defense mechanisms discussed in the previous section. Adversarial training involves training a model on adversarial examples to improve its robustness, while robustness certification provides a guarantee that a model will behave correctly even under adversarial attacks. In this subsection, we will delve into the specifics of these techniques, including adversarial example generation, robust optimization, and certified defenses, and explore how they can be used to enhance the security and reliability of computer vision systems.

Adversarial example generation is a critical component of adversarial training, as it allows models to learn from examples that are specifically designed to mislead them \cite{ref62}. One popular method for generating adversarial examples is the Fast Gradient Sign Method (FGSM) \cite{ref110}, which uses the gradient of the model's loss function to create adversarial examples. Another method is the Projected Gradient Descent (PGD) attack \cite{ref111}, which uses an iterative approach to generate adversarial examples. These methods can be used to generate a wide range of adversarial examples, allowing models to learn to defend against different types of attacks.

Robust optimization is another key aspect of adversarial training, as it involves modifying the model's objective function to include a term that encourages robustness \cite{ref112}. One popular method for robust optimization is adversarial training with a minimax objective \cite{ref113}, which involves minimizing the model's loss function while maximizing the loss function of an adversary. Another method is robust optimization using a regularization term \cite{ref114}, which adds a penalty term to the model's objective function to encourage robustness. These methods can be used to improve the robustness of models to adversarial attacks, and can be combined with adversarial example generation to create a comprehensive defense strategy.

Certified defenses provide a guarantee that a model will behave correctly even under adversarial attacks \cite{ref115}, and are an essential component of any defense strategy. One popular method for certified defenses is Interval Bound Propagation (IBP) \cite{ref116}, which involves propagating interval bounds through the model to provide a guarantee of robustness. Another method is Certified Adversarial Training (CAT) \cite{ref117}, which involves training a model to be robust to adversarial attacks while providing a certificate of robustness. These methods can be used to provide a guarantee of robustness, and can be combined with adversarial training and robust optimization to create a comprehensive defense strategy.

In addition to these techniques, there are several other methods that have been proposed for adversarial training and robustness certification, including the use of generative models to generate adversarial examples \cite{ref118}, and the use of reinforcement learning to train models to be robust to adversarial attacks \cite{ref119}. These methods can be used to improve the robustness of models to adversarial attacks, and can be combined with other defense strategies to create a comprehensive defense system.

Despite the many advances that have been made in adversarial training and robustness certification, there are still several challenges that remain to be addressed, including the trade-off between robustness and accuracy \cite{ref120}. Many methods for adversarial training and robustness certification involve sacrificing some accuracy in order to improve robustness, and finding a balance between these two competing goals is an ongoing challenge. Another challenge is the scalability of these methods to large models and datasets \cite{ref121}, as many methods for adversarial training and robustness certification are computationally expensive and may not be feasible for large models and datasets.

To address these challenges, researchers have proposed several new methods, including the use of label smoothing and logit squeezing \cite{ref122}, and the use of game-theoretic approaches \cite{ref123}. These methods can be used to improve the robustness of models to adversarial attacks, and can be combined with other defense strategies to create a comprehensive defense system. Additionally, the use of adaptive continuous adversarial training (ACAT) \cite{ref124} and vulnerability-aware and curiosity-driven adversarial training (VCAT) \cite{ref125} have been proposed to enhance the robustness of models against adversarial attacks.

In conclusion, adversarial training and robustness certification are two crucial defense mechanisms against visual adversarial attacks, and are essential components of any comprehensive defense strategy. By combining these techniques with other defense strategies, such as input preprocessing and detection methods, we can create a robust and comprehensive defense system that can protect against a wide range of adversarial attacks. The development of more effective and efficient methods for adversarial training and robustness certification is critical to ensuring the security and reliability of computer vision systems, and is an area of ongoing research and development \cite{ref126}. As we move forward, it is essential that we continue to develop and refine these methods, and to explore new and innovative approaches to defending against visual adversarial attacks.

\subsection{Input Preprocessing and Detection Methods}

Input preprocessing and detection methods have emerged as a crucial defense mechanism against visual adversarial attacks \cite{ref127}. These methods aim to identify and eliminate adversarial perturbations from the input data, thereby preventing the attack from succeeding. One approach is to use image reconstruction techniques, such as image compression \cite{ref127} or super-resolution \cite{ref127}, to remove adversarial noise from the input images. Another approach is to utilize digital image processing techniques, such as spatial filtering \cite{ref128} or thresholding \cite{ref128}, to detect and remove adversarial perturbations.

Detection-based methods, on the other hand, focus on identifying adversarial examples by analyzing the input data \cite{ref129}. One technique is to use a secondary model, such as a detector network \cite{ref129}, to detect adversarial examples. Another approach is to analyze the input data's statistical properties, such as the distribution of pixel values \cite{ref130}, to identify adversarial examples.

Preprocessing-based methods can be further divided into two categories: passive and active. Passive methods, such as input validation \cite{ref131} or data normalization \cite{ref131}, aim to prevent adversarial examples from being input into the system. Active methods, such as adversarial training \cite{ref106} or input purification \cite{ref132}, aim to detect and eliminate adversarial perturbations from the input data.

Several papers have demonstrated the effectiveness of input preprocessing and detection methods in defending against visual adversarial attacks \cite{ref128,ref127}. For example, \cite{ref128} proposed a defense method based on digital image processing, which can effectively defend against adversarial attacks. \cite{ref127} proposed an image reconstruction-based defense method, which can remove adversarial noise from input images.

However, input preprocessing and detection methods also have some limitations \cite{ref133}. For example, they may not be effective against all types of adversarial attacks, such as attacks that use multiple types of perturbations \cite{ref133}. Additionally, these methods may require significant computational resources, which can be a challenge for real-time applications \cite{ref134}.

In conclusion, input preprocessing and detection methods are a crucial defense mechanism against visual adversarial attacks \cite{ref6}. These methods can be used to identify and eliminate adversarial perturbations from the input data, thereby preventing the attack from succeeding. While these methods have some limitations, they have been shown to be effective in defending against various types of adversarial attacks \cite{ref128,ref127}. Further research is needed to improve the effectiveness and efficiency of these methods, as well as to develop new methods that can defend against emerging types of adversarial attacks \cite{ref6}.

To further enhance the robustness of deep neural networks against adversarial attacks, some papers have proposed using basis functions transformations \cite{ref51} or PCA \cite{ref135} to defend against adversarial attacks. These methods can be used to remove adversarial noise from the input data, thereby preventing the attack from succeeding. Additionally, some papers have proposed using deep latent defense \cite{ref136} or robust adversarial defense by tensor factorization \cite{ref137} to defend against adversarial attacks. These methods can be used to detect and eliminate adversarial perturbations from the input data, thereby preventing the attack from succeeding.

Moreover, some papers have proposed using reinforcement learning \cite{ref105} or self-supervised learning \cite{ref138} to improve the robustness of deep neural networks against adversarial attacks. These methods can be used to improve the robustness of deep neural networks against adversarial attacks, thereby preventing the attack from succeeding. Furthermore, some papers have proposed using deep learning-based methods \cite{ref139} or computer vision-based methods \cite{ref6} to defend against adversarial attacks. These methods can be used to improve the robustness of deep neural networks against adversarial attacks, thereby preventing the attack from succeeding.

It is also worth noting that some papers have proposed using other types of defense mechanisms, such as certified defenses \cite{ref39} or robust optimization methods \cite{ref106}. These methods can be used to improve the robustness of deep neural networks against adversarial attacks, thereby preventing the attack from succeeding. Overall, input preprocessing and detection methods are a crucial defense mechanism against visual adversarial attacks \cite{ref6}. These methods can be used to identify and eliminate adversarial perturbations from the input data, thereby preventing the attack from succeeding. While these methods have some limitations, they have been shown to be effective in defending against various types of adversarial attacks \cite{ref128,ref127}. Further research is needed to improve the effectiveness and efficiency of these methods, as well as to develop new methods that can defend against emerging types of adversarial attacks \cite{ref6}.

\section{Evaluation Metrics and Benchmarking for Adversarial Robustness}

\subsection{Introduction to Evaluation Metrics}

Evaluation metrics play a crucial role in assessing the robustness of computer vision models against adversarial attacks \cite{ref140}. The importance of these metrics cannot be overstated, as they provide a standardized way to measure the performance of models under various attack scenarios. In the context of adversarial robustness, evaluation metrics help researchers and practitioners understand the vulnerabilities of their models and identify areas for improvement \cite{ref16}.

One of the primary challenges in evaluating adversarial robustness is the lack of a unified framework for measuring robustness \cite{ref141}. Different metrics may yield varying results, making it difficult to compare the robustness of different models \cite{ref142}. To address this issue, researchers have proposed various evaluation metrics, including the use of \(L_p\) norms, such as \(L_0\), \(L_2\), and \(L_\infty\), to measure the size of perturbations \cite{ref141}.

In addition to \(L_p\) norms, other metrics, such as the Wasserstein distance, have been proposed to measure the robustness of models \cite{ref143}. The Wasserstein distance takes into account the geometric properties of the input data, providing a more nuanced measure of robustness \cite{ref144}. However, the choice of evaluation metric ultimately depends on the specific use case and the type of attack being considered \cite{ref21}.

The development of evaluation metrics is an active area of research, with new metrics and methods being proposed regularly \cite{ref145}. For example, researchers have proposed the use of ensemble methods, which combine the predictions of multiple models to improve robustness \cite{ref146}. Other approaches, such as adversarial training, involve training models on adversarial examples to improve their robustness \cite{ref147}.

In conclusion, evaluation metrics play a vital role in assessing the robustness of computer vision models against adversarial attacks \cite{ref148}. The choice of metric depends on the specific use case and the type of attack being considered \cite{ref6}. Researchers and practitioners must carefully consider the strengths and limitations of different metrics and use them in conjunction with other methods, such as adversarial training and ensemble methods, to develop more robust models \cite{ref149}. By doing so, we can improve the reliability and security of computer vision systems in a wide range of applications \cite{ref150}.

The development of new evaluation metrics and methods will be crucial for improving the robustness of computer vision models \cite{ref143}. Researchers and practitioners must continue to develop and refine evaluation metrics to stay ahead of the evolving landscape of adversarial attacks \cite{ref151}. The importance of evaluation metrics in assessing the robustness of computer vision models against adversarial attacks will only continue to grow as the use of computer vision systems becomes more widespread \cite{ref152}.

Therefore, it is essential to continue developing and refining evaluation metrics to improve the robustness of computer vision models \cite{ref153}. The use of evaluation metrics, in conjunction with other methods, such as adversarial training and ensemble methods, can help improve the reliability and security of computer vision systems \cite{ref154}. By doing so, we can develop more robust models that are better equipped to withstand the evolving landscape of adversarial attacks \cite{ref155}.

In the future, the development of new evaluation metrics and methods will be crucial for improving the robustness of computer vision models \cite{ref90}. Researchers and practitioners must continue to develop and refine evaluation metrics to stay ahead of the evolving landscape of adversarial attacks \cite{ref20}. The importance of evaluation metrics in assessing the robustness of computer vision models against adversarial attacks will only continue to grow as the use of computer vision systems becomes more widespread \cite{ref106}.

In conclusion, evaluation metrics play a vital role in assessing the robustness of computer vision models against adversarial attacks \cite{ref156}. The choice of metric depends on the specific use case and the type of attack being considered \cite{ref157}. Researchers and practitioners must carefully consider the strengths and limitations of different metrics and use them in conjunction with other methods, such as adversarial training and ensemble methods, to develop more robust models \cite{ref158}. By doing so, we can improve the reliability and security of computer vision systems in a wide range of applications \cite{ref159}.

The development of new evaluation metrics and methods is crucial for improving the robustness of models and identifying potential vulnerabilities \cite{ref160}. Researchers and practitioners must continue to develop and refine evaluation metrics to stay ahead of the evolving landscape of adversarial attacks \cite{ref161}. The importance of evaluation metrics in assessing the robustness of computer vision models against adversarial attacks cannot be overstated \cite{ref162}. The use of evaluation metrics can help identify areas for improvement in model robustness \cite{ref163}. By analyzing the results of evaluation metrics, researchers and practitioners can develop more effective defense strategies and improve the overall robustness of computer vision models \cite{ref164}.

Therefore, it is essential to continue developing and refining evaluation metrics to improve the robustness of computer vision models \cite{ref165}. The use of evaluation metrics, in conjunction with other methods, such as adversarial training and ensemble methods, can help improve the reliability and security of computer vision systems \cite{ref21}. By doing so, we can develop more robust models that are better equipped to withstand the evolving landscape of adversarial attacks \cite{ref166}.

In the future, the development of new evaluation metrics and methods will be crucial for improving the robustness of computer vision models \cite{ref90}. Researchers and practitioners must continue to develop and refine evaluation metrics to stay ahead of the evolving landscape of adversarial attacks \cite{ref20}. The importance of evaluation metrics in assessing the robustness of computer vision models against adversarial attacks will only continue to grow as the use of computer vision systems becomes more widespread \cite{ref106}.

In conclusion, evaluation metrics play a vital role in assessing the robustness of computer vision models against adversarial attacks \cite{ref148}. The choice of metric depends on the specific use case and the type of attack being considered \cite{ref6}. Researchers and practitioners must carefully consider the strengths and limitations of different metrics and use them in conjunction with other methods, such as adversarial training and ensemble methods, to develop more robust models \cite{ref149}. By doing so, we can improve the reliability and security of computer vision systems in a wide range of applications \cite{ref150}.

The development of new evaluation metrics and methods is crucial for improving the robustness of models and identifying potential vulnerabilities \cite{ref167}. Researchers and practitioners must continue to develop and refine evaluation metrics to stay ahead of the evolving landscape of adversarial attacks \cite{ref168}. The importance of evaluation metrics in assessing the robustness of computer vision models against adversarial attacks cannot be overstated \cite{ref169}. The use of evaluation metrics can help identify areas for improvement in model robustness \cite{ref170}. By analyzing the results of evaluation metrics, researchers and practitioners can develop more effective defense strategies and improve the overall robustness of computer vision models \cite{ref171}.

Therefore, it is essential to continue developing and refining evaluation metrics to improve the robustness of computer vision models \cite{ref153}. The use of evaluation metrics, in conjunction with other methods, such as adversarial training and ensemble methods, can help improve the reliability and security of computer vision systems \cite{ref154}. By doing so, we can develop more robust models that are better equipped to withstand the evolving landscape of adversarial attacks \cite{ref155}.

In the future, the development of new evaluation metrics and methods will be crucial for improving the robustness of computer vision models \cite{ref143}. Researchers and practitioners must continue to develop and refine evaluation metrics to stay ahead of the evolving landscape of adversarial attacks \cite{ref151}. The importance of evaluation metrics in assessing the robustness of computer vision models against adversarial attacks will only continue to grow as the use of computer vision systems becomes more widespread \cite{ref152}.

In conclusion, evaluation metrics play a vital role in assessing the robustness of computer vision models against adversarial attacks \cite{ref156}. The choice of metric depends on the specific use case and the type of attack being considered \cite{ref157}. Researchers and practitioners must carefully consider the strengths and limitations of different metrics and use them in conjunction with other methods, such as adversarial training and ensemble methods, to develop more robust models \cite{ref158}. By doing so, we can improve the reliability and security of computer vision systems in a wide range of applications \cite{ref159}.

The development of new evaluation metrics and methods is crucial for improving the robustness of models and identifying potential vulnerabilities \cite{ref160}. Researchers and practitioners must continue to develop and refine evaluation metrics to stay ahead of the evolving landscape of adversarial attacks \cite{ref161}. The importance of evaluation metrics in assessing the robustness of computer vision models against adversarial attacks cannot be overstated \cite{ref162}. The use of evaluation metrics can help identify areas for improvement in model robustness \cite{ref163}. By analyzing the results of evaluation metrics, researchers and practitioners can develop more effective defense strategies and improve the overall robustness of computer vision models \cite{ref164}.

Therefore, it is essential to continue developing and refining evaluation metrics to improve the robustness of computer vision models \cite{ref165}. The use of evaluation metrics, in conjunction with other methods, such as adversarial training and ensemble methods, can help improve the reliability and security of computer vision systems \cite{ref21}. By doing so, we can develop more robust models that are better equipped to withstand the evolving landscape of adversarial attacks \cite{ref166}.

In the future, the development of new evaluation metrics and methods will be crucial for improving the robustness of computer vision models \cite{ref90}. Researchers and practitioners must continue to develop and refine evaluation metrics to stay ahead of the evolving landscape of adversarial attacks \cite{ref20}. The importance of evaluation metrics in assessing the robustness of computer vision models against adversarial attacks will only continue to grow as the use of computer vision systems becomes more widespread \cite{ref106}.

In conclusion, evaluation metrics play a vital role in assessing the robustness of computer vision models against adversarial attacks \cite{ref148}. The choice of metric depends on the specific use case and the type of attack being considered \cite{ref6}. Researchers and practitioners must carefully consider the strengths and limitations of different metrics and use them in conjunction with other methods, such as adversarial training and ensemble methods, to develop more robust models \cite{ref149}. By doing so, we can improve the reliability and security of computer vision systems in a wide range of applications \cite{ref150}.

The development of new evaluation metrics and methods is crucial for improving the robustness of models and identifying potential vulnerabilities \cite{ref167}. Researchers and practitioners must continue to develop and refine evaluation metrics to stay ahead of the evolving landscape of adversarial attacks \cite{ref168}. The importance of evaluation metrics in assessing the robustness of computer vision models against adversarial attacks cannot be overstated \cite{ref169}. The use of evaluation metrics can help identify areas for improvement in model robustness \cite{ref170}. By analyzing the results of evaluation metrics, researchers and practitioners can develop more effective defense strategies and improve the overall robustness of computer vision models \cite{ref171}.

Therefore, it is essential to continue developing and refining evaluation metrics to improve the robustness of computer vision models \cite{ref153}. The use of evaluation metrics, in conjunction with other methods, such as adversarial training and ensemble methods, can help improve the reliability and security of computer vision systems \cite{ref154}. By doing so, we can develop more robust models that are better equipped to withstand the evolving landscape of adversarial attacks \cite{ref155}.

In the future, the development of new evaluation metrics and methods will be crucial for improving the robustness of computer vision models \cite{ref143}. Researchers and practitioners must continue to develop and refine evaluation metrics to stay ahead of the evolving landscape of adversarial attacks \cite{ref151}. The importance of evaluation metrics in assessing the robustness of computer vision models against adversarial attacks will only continue to grow as the use of computer vision systems becomes more widespread \cite{ref152}.

In conclusion, evaluation metrics play a vital role in assessing the robustness of computer vision models against adversarial attacks \cite{ref156}. The choice of metric depends on the specific use case and the type of attack being considered \cite{ref157}. Researchers and practitioners must carefully consider the strengths and limitations of different metrics and use them in conjunction with other methods, such as adversarial training and ensemble methods, to develop more robust models \cite{ref158}. By doing so, we can improve the reliability and security of computer vision systems in a wide range of applications \cite{ref159}.

The development of new evaluation metrics and methods is crucial for improving the robustness of models and identifying potential vulnerabilities \cite{ref160}. Researchers and practitioners must continue to develop and refine evaluation metrics to stay ahead of the evolving landscape of adversarial attacks \cite{ref161}. The importance of evaluation metrics in assessing the robustness of computer vision models against adversarial attacks cannot be overstated \cite{ref162}. The use of evaluation metrics can help identify areas for improvement in model robustness \cite{ref163}. By analyzing the results of evaluation metrics, researchers and practitioners can develop more effective defense strategies and improve the overall robustness of computer vision models \cite{ref164}.

Therefore, it is essential to continue developing and refining evaluation metrics to improve the robustness of computer vision models \cite{ref165}. The use of evaluation metrics, in conjunction with other methods, such as adversarial training and ensemble methods, can help improve the reliability and security of computer vision systems \cite{ref21}. By doing so, we can develop more robust models that are better equipped to withstand the evolving landscape of adversarial attacks \cite{ref166}.

In the future, the development of new evaluation metrics and methods will be crucial for improving the robustness of computer vision models \cite{ref90}. Researchers and practitioners must continue to develop and refine evaluation metrics to stay ahead of the evolving landscape of adversarial attacks \cite{ref20}. The importance of evaluation metrics in assessing the robustness of computer vision models against adversarial attacks will only continue to grow as the use of computer vision systems becomes more widespread \cite{ref106}.

In conclusion, evaluation metrics play a vital role in assessing the robustness of computer vision models against adversarial attacks \cite{ref148}. The choice of metric depends on the specific use case and the type of attack being considered \cite{ref6}. Researchers and practitioners must carefully consider the strengths and limitations of different metrics and use them in conjunction with other methods, such as adversarial training and ensemble methods, to develop more robust models \cite{ref149}. By doing so, we can improve the reliability and security of computer vision systems in a wide range of applications \cite{ref150}.

The development of new evaluation metrics and methods is crucial for improving the robustness of models and identifying potential vulnerabilities \cite{ref167}. Researchers and practitioners must continue to develop and refine evaluation metrics to stay ahead of the evolving landscape of adversarial attacks \cite{ref168}. The importance of evaluation metrics in assessing the robustness of computer vision models against adversarial attacks cannot be overstated \cite{ref169}. The use of evaluation metrics can help identify areas for improvement in model robustness \cite{ref170}. By analyzing the results of evaluation metrics, researchers and practitioners can develop more effective defense strategies and improve the overall robustness of computer vision models \cite{ref171}.

Therefore, it is essential to continue developing and refining evaluation metrics to improve the robustness of computer vision models \cite{ref153}. The use of evaluation metrics, in conjunction with other methods, such as adversarial training and ensemble methods, can help improve the reliability and security of computer vision systems \cite{ref154}. By doing so, we can develop more robust models that are better equipped to withstand the evolving landscape of adversarial attacks \cite{ref155}.

In the future, the development of new evaluation metrics and methods will be crucial for improving the robustness of computer vision models \cite{ref143}. Researchers and practitioners must continue to develop and refine evaluation metrics to stay ahead of the evolving landscape of adversarial attacks \cite{ref151}. The importance of evaluation metrics in assessing the robustness of computer vision models against adversarial attacks will only continue to grow as the use of computer vision systems becomes more widespread \cite{ref152}.

In conclusion, evaluation metrics play a vital role in assessing the robustness of computer vision models against adversarial attacks \cite{ref156}. The choice of metric depends on the specific use case and the type of attack being considered \cite{ref157}. Researchers and practitioners must carefully consider the strengths and limitations of different metrics and use them in conjunction with other methods, such as adversarial training and ensemble methods, to develop more robust models \cite{ref158}. By doing so, we can improve the reliability and security of computer vision systems in a wide range of applications \cite{ref159}.

The development of new evaluation metrics and methods is crucial for improving the robustness of models and identifying potential vulnerabilities \cite{ref160}. Researchers and practitioners must continue to develop and refine evaluation metrics to stay ahead of the evolving landscape of adversarial attacks \cite{ref161}. The importance of evaluation metrics in assessing the robustness of computer vision models against adversarial attacks cannot be overstated \cite{ref162}. The use of evaluation metrics can help identify areas for improvement in model robustness \cite{ref163}. By analyzing the results of evaluation metrics, researchers and practitioners can develop more effective defense strategies and improve the overall robustness of computer vision models \cite{ref164}.

Therefore, it is essential to continue developing and refining evaluation metrics to improve the robustness of computer vision models \cite{ref165}. The use of evaluation metrics, in conjunction with other methods, such as adversarial training and ensemble methods, can help improve the reliability and security of computer vision systems \cite{ref21}. By doing so, we can develop more robust models that are better equipped to withstand the evolving landscape of adversarial attacks \cite{ref166}.

In the future, the development of new evaluation metrics and methods will be crucial for improving the robustness of computer vision models \cite{ref90}. Researchers and practitioners must continue to develop and refine evaluation metrics to stay ahead of the evolving landscape of adversarial attacks \cite{ref20}. The importance of evaluation metrics in assessing the robustness of computer vision models against adversarial attacks will only continue to grow as the use of computer vision systems becomes more widespread \cite{ref106}.

In conclusion, evaluation metrics play a vital role in assessing the robustness of computer vision models against adversarial attacks \cite{ref148}. The choice of metric depends on the specific use case and the type of attack being considered \cite{ref6}. Researchers and practitioners must carefully consider the strengths and limitations of different metrics and use them in conjunction with other methods, such as adversarial training and ensemble methods, to develop more robust models \cite{ref149}. By doing so, we can improve the reliability and security of computer vision systems in a wide range of applications \cite{ref150}.

The development of new evaluation metrics and methods is crucial for improving the robustness of models and identifying potential vulnerabilities \cite{ref167}. Researchers and practitioners must continue to develop and refine evaluation metrics to stay ahead of the evolving landscape of adversarial attacks \cite{ref168}. The importance of evaluation metrics in assessing the robustness of computer vision models against adversarial attacks cannot be overstated \cite{ref169}. The use of evaluation metrics can help identify areas for improvement in model robustness \cite{ref170}. By analyzing the results of evaluation metrics, researchers and practitioners can develop more effective defense strategies and improve the overall robustness of computer vision models \cite{ref171}.

Therefore, it is essential to continue developing and refining evaluation metrics to improve the robustness of computer vision models \cite{ref153}. The use of evaluation metrics, in conjunction with other methods, such as adversarial training and ensemble methods, can help improve the reliability and security of computer vision systems \cite{ref154}. By doing so, we can develop more robust models that are better equipped to withstand the evolving landscape of adversarial attacks \cite{ref155}.

In the future, the development of new evaluation metrics and methods will be crucial for improving the robustness of computer vision models \cite{ref143}. Researchers and practitioners must continue to develop and refine evaluation metrics to stay ahead of the evolving landscape of adversarial attacks \cite{ref151}. The importance of evaluation metrics in assessing the robustness of computer vision models against adversarial attacks will only continue to grow as the use of computer vision systems becomes more widespread \cite{ref152}.

In conclusion, evaluation metrics play a vital role in assessing the robustness of computer vision models against adversarial attacks \cite{ref156}. The choice of metric depends on the specific use case and the type of attack being considered \cite{ref157}. Researchers and practitioners must carefully consider the strengths and limitations of different metrics and use them in conjunction with other methods, such as adversarial training and ensemble methods, to develop more robust models \cite{ref158}. By doing so, we can improve the reliability and security of computer vision systems in a wide range of applications \cite{ref159}.

The development of new evaluation metrics and methods is crucial for improving the robustness of models and identifying potential vulnerabilities \cite{ref160}. Researchers and practitioners must continue to develop and refine evaluation metrics to stay ahead of the evolving landscape of adversarial attacks \cite{ref161}. The importance of evaluation metrics in assessing the robustness of computer vision models against adversarial attacks cannot be overstated \cite{ref162}. The use of evaluation metrics can help identify areas for improvement in model robustness \cite{ref163}. By analyzing the results of evaluation metrics, researchers and practitioners can develop more effective defense strategies and improve the overall robustness of computer vision models \cite{ref164}.

Therefore, it is essential to continue developing and refining evaluation metrics to improve the robustness of computer vision models \cite{ref165}. The use of evaluation metrics, in conjunction with other methods, such as adversarial training and ensemble methods, can help improve the reliability and security of computer vision systems \cite{ref21}. By doing so, we can develop more robust models that are better equipped to withstand the evolving landscape of adversarial attacks \cite{ref166}.

In the future, the development of new evaluation metrics and methods will be crucial for improving the robustness of computer vision models \cite{ref90}. Researchers and practitioners must continue to develop and refine evaluation metrics to stay ahead of the evolving landscape of adversarial attacks \cite{ref20}. The importance of evaluation metrics in assessing the robustness of computer vision models against adversarial attacks will only continue to grow as the use of computer vision systems becomes more widespread \cite{ref106}.

In conclusion, evaluation metrics play a vital role in assessing the robustness of computer vision models against adversarial attacks \cite{ref148}. The choice of metric depends on the specific use case and the type of attack being considered \cite{ref6}. Researchers and practitioners must carefully consider the strengths and limitations of different metrics and use them in conjunction with other methods, such as adversarial training and ensemble methods, to develop more robust models \cite{ref149}. By doing so, we can improve the reliability and security of computer vision systems in a wide range of applications \cite{ref150}.

The development of new evaluation metrics and methods is crucial for improving the robustness of models and identifying potential vulnerabilities \cite{ref167}. Researchers and practitioners must continue to develop and refine evaluation metrics to stay ahead of the evolving landscape of adversarial attacks \cite{ref168}. The importance of evaluation metrics in assessing the robustness of computer vision models against adversarial attacks cannot be overstated \cite{ref169}. The use of evaluation metrics can help identify areas for improvement in model robustness \cite{ref170}. By analyzing the results of evaluation metrics, researchers and practitioners can develop more effective defense strategies and improve the overall robustness of computer vision models \cite{ref171}.

Therefore, it is essential to continue developing and refining evaluation metrics to improve the robustness of computer vision models \cite{ref153}. The use of evaluation metrics, in conjunction with other methods, such as adversarial training and ensemble methods, can help improve the reliability and security of computer vision systems \cite{ref154}. By doing so, we can develop more robust models that are better equipped to withstand the evolving landscape of adversarial attacks \cite{ref155}.

In the future, the development of new evaluation metrics and methods will be crucial for improving the robustness of computer vision models \cite{ref143}. Researchers and practitioners must continue to develop and refine evaluation metrics to stay ahead of the evolving landscape of adversarial attacks \cite{ref151}. The importance of evaluation metrics in assessing the robustness of computer vision models against adversarial attacks will only continue to grow as the use of computer vision systems becomes more widespread \cite{ref152}.

In conclusion, evaluation metrics play a vital role in assessing the robustness of computer vision models against adversarial attacks \cite{ref156}. The choice of metric depends on the specific use case and the type of attack being considered \cite{ref157}. Researchers and practitioners must carefully consider the strengths and limitations of different metrics and use them in conjunction with other methods, such as adversarial training and ensemble methods, to develop more robust models \cite{ref158}. By doing so, we can improve the reliability and security of computer vision systems in a wide range of applications \cite{ref159}.

The development of new evaluation metrics and methods is crucial for improving the robustness of models and identifying potential vulnerabilities \cite{ref160}. Researchers and practitioners must continue to develop and refine evaluation metrics to stay ahead of the evolving landscape of adversarial attacks \cite{ref161}. The importance of evaluation metrics in assessing the robustness of computer vision models against adversarial attacks cannot be overstated \cite{ref162}. The use of evaluation metrics can help identify areas for improvement in model robustness \cite{ref163}. By analyzing the results of evaluation metrics, researchers and practitioners can develop more effective defense strategies and improve the overall robustness of computer vision models \cite{ref164}.

Therefore, it is essential to continue developing and refining evaluation metrics to improve the robustness of computer vision models \cite{ref165}. The use of evaluation metrics, in conjunction with other methods, such as adversarial training and ensemble methods, can help improve the reliability and security of computer vision systems \cite{ref21}. By doing so, we can develop more robust models that are better equipped to withstand the evolving landscape of adversarial attacks \cite{ref166}.

In the future, the development of new evaluation metrics and methods will be crucial for improving the robustness of computer vision models \cite{ref90}. Researchers and practitioners must continue to develop and refine evaluation metrics to stay ahead of the evolving landscape of adversarial attacks \cite{ref20}. The importance of evaluation metrics in assessing the robustness of computer vision models against adversarial attacks will only continue to grow as the use of computer vision systems becomes more widespread \cite{ref106}.

In conclusion, evaluation metrics play a vital role in assessing the robustness of computer vision models against adversarial attacks \cite{ref148}. The choice of metric depends on the specific use case and the type of attack being considered \cite{ref6}. Researchers and practitioners must carefully consider the strengths and limitations of different metrics and use them in conjunction with other methods, such as adversarial training and ensemble methods, to develop more robust models \cite{ref149}. By doing so, we can improve the reliability and security of computer vision systems in a wide range of applications \cite{ref150}.

The development of new evaluation metrics and methods is crucial for improving the robustness of models and identifying potential vulnerabilities \cite{ref167}. Researchers and practitioners must continue to develop and refine evaluation metrics to stay ahead of the evolving landscape of adversarial attacks \cite{ref168}. The importance of evaluation metrics in assessing the robustness of computer vision

\subsection{Metrics for Adversarial Robustness}

Metrics for Adversarial Robustness are crucial in evaluating the effectiveness of defense mechanisms against adversarial attacks. Various metrics have been proposed to assess the robustness of deep neural networks, including L0, L\(\infty\), and L2 norms. The L0 norm measures the number of pixels that have been modified, while the L\(\infty\) norm measures the maximum change in any pixel value \cite{ref142}. The L2 norm, on the other hand, measures the Euclidean distance between the original and perturbed images \cite{ref172}.

In addition to these norms, other metrics have been proposed to evaluate the effectiveness of defense mechanisms. For example, the robustness metric proposed in \cite{ref100} assesses the robustness of a model by measuring the minimum perturbation required to misclassify an image. This metric takes into account both the L0 and L\(\infty\) norms, providing a more comprehensive evaluation of a model's robustness.

Another important metric is the attack success rate (ASR), which measures the percentage of successful attacks on a model \cite{ref173}. The ASR metric is widely used to evaluate the effectiveness of defense mechanisms, as it provides a clear indication of a model's vulnerability to adversarial attacks.

Furthermore, metrics such as the mean squared error (MSE) and peak signal-to-noise ratio (PSNR) have been used to evaluate the quality of adversarial examples \cite{ref174}. These metrics provide a measure of the similarity between the original and perturbed images, allowing for the evaluation of the perceptual quality of adversarial examples.

The development of effective metrics for evaluating adversarial robustness is an ongoing area of research, with many new metrics and evaluation frameworks being proposed in recent years. As the field of adversarial robustness continues to evolve, the development of effective metrics will remain a critical area of research, enabling the evaluation and improvement of defense mechanisms against adversarial attacks \cite{ref175}.

In conclusion, the evaluation of adversarial robustness is a critical area of research, requiring the development of effective metrics and evaluation frameworks. The metrics mentioned above, including L0, L\(\infty\), and L2 norms, as well as more sophisticated metrics that take into account the model's internal workings and the geometry of the input data, provide a comprehensive evaluation of a model's robustness. The development of effective metrics for evaluating adversarial robustness will remain a critical area of research, enabling the evaluation and improvement of defense mechanisms against adversarial attacks.

This subsection provides a foundation for understanding the importance of metrics in evaluating adversarial robustness, which will be further explored in the following subsection on benchmarking and evaluation frameworks. The development of effective metrics and evaluation frameworks is essential for advancing the field of adversarial robustness and ensuring the security and reliability of deep neural networks in real-world applications \cite{ref156}.

As the field of adversarial robustness continues to evolve, the development of effective metrics will play a critical role in enabling the evaluation and improvement of defense mechanisms against adversarial attacks \cite{ref176}. The importance of evaluating adversarial robustness cannot be overstated, as deep neural networks are increasingly used in safety-critical applications, and the need to evaluate and improve their robustness against adversarial attacks becomes more pressing \cite{ref162}.

The use of effective metrics and evaluation frameworks will enable researchers to develop more robust models, identify areas for improvement, and advance the field of adversarial robustness \cite{ref171}. The development of effective metrics for evaluating adversarial robustness is an ongoing area of research, with many new metrics and evaluation frameworks being proposed in recent years \cite{ref144}.

In summary, the evaluation of adversarial robustness is a complex and multifaceted problem, requiring the development of effective metrics and evaluation frameworks. The metrics mentioned above provide a comprehensive evaluation of a model's robustness, and the development of effective metrics will remain a critical area of research, enabling the evaluation and improvement of defense mechanisms against adversarial attacks \cite{ref177}.

The importance of evaluating adversarial robustness cannot be overstated, and the development of effective metrics and evaluation frameworks is essential for advancing the field of adversarial robustness \cite{ref178}. The use of effective metrics and evaluation frameworks will enable researchers to develop more robust models, identify areas for improvement, and advance the field of adversarial robustness \cite{ref179}.

In conclusion, the evaluation of adversarial robustness is a critical area of research, requiring the development of effective metrics and evaluation frameworks. The development of effective metrics for evaluating adversarial robustness will remain a critical area of research, enabling the evaluation and improvement of defense mechanisms against adversarial attacks \cite{ref180}.

The development of effective metrics for evaluating adversarial robustness is an ongoing area of research, with many new metrics and evaluation frameworks being proposed in recent years \cite{ref145}. The importance of evaluating adversarial robustness cannot be overstated, and the development of effective metrics and evaluation frameworks is essential for advancing the field of adversarial robustness \cite{ref181}.

In summary, the evaluation of adversarial robustness is a complex and multifaceted problem, requiring the development of effective metrics and evaluation frameworks. The metrics mentioned above provide a comprehensive evaluation of a model's robustness, and the development of effective metrics will remain a critical area of research, enabling the evaluation and improvement of defense mechanisms against adversarial attacks \cite{ref182}.

The importance of evaluating adversarial robustness cannot be overstated, and the development of effective metrics and evaluation frameworks is essential for advancing the field of adversarial robustness \cite{ref140}. The use of effective metrics and evaluation frameworks will enable researchers to develop more robust models, identify areas for improvement, and advance the field of adversarial robustness \cite{ref183}.

In conclusion, the evaluation of adversarial robustness is a critical area of research, requiring the development of effective metrics and evaluation frameworks. The development of effective metrics for evaluating adversarial robustness will remain a critical area of research, enabling the evaluation and improvement of defense mechanisms against adversarial attacks \cite{ref184}.

The development of effective metrics for evaluating adversarial robustness is an ongoing area of research, with many new metrics and evaluation frameworks being proposed in recent years \cite{ref185}. The importance of evaluating adversarial robustness cannot be overstated, and the development of effective metrics and evaluation frameworks is essential for advancing the field of adversarial robustness \cite{ref186}.

In summary, the evaluation of adversarial robustness is a complex and multifaceted problem, requiring the development of effective metrics and evaluation frameworks. The metrics mentioned above provide a comprehensive evaluation of a model's robustness, and the development of effective metrics will remain a critical area of research, enabling the evaluation and improvement of defense mechanisms against adversarial attacks \cite{ref187}.

The importance of evaluating adversarial robustness cannot be overstated, and the development of effective metrics and evaluation frameworks is essential for advancing the field of adversarial robustness \cite{ref188}. The use of effective metrics and evaluation frameworks will enable researchers to develop more robust models, identify areas for improvement, and advance the field of adversarial robustness \cite{ref189}.

In conclusion, the evaluation of adversarial robustness is a critical area of research, requiring the development of effective metrics and evaluation frameworks. The development of effective metrics for evaluating adversarial robustness will remain a critical area of research, enabling the evaluation and improvement of defense mechanisms against adversarial attacks \cite{ref190}.

The development of effective metrics for evaluating adversarial robustness is an ongoing area of research, with many new metrics and evaluation frameworks being proposed in recent years \cite{ref144}. The importance of evaluating adversarial robustness cannot be overstated, and the development of effective metrics and evaluation frameworks is essential for advancing the field of adversarial robustness \cite{ref177}.

In summary, the evaluation of adversarial robustness is a complex and multifaceted problem, requiring the development of effective metrics and evaluation frameworks. The metrics mentioned above provide a comprehensive evaluation of a model's robustness, and the development of effective metrics will remain a critical area of research, enabling the evaluation and improvement of defense mechanisms against adversarial attacks \cite{ref178}.

The importance of evaluating adversarial robustness cannot be overstated, and the development of effective metrics and evaluation frameworks is essential for advancing the field of adversarial robustness \cite{ref179}. The use of effective metrics and evaluation frameworks will enable researchers to develop more robust models, identify areas for improvement, and advance the field of adversarial robustness \cite{ref180}.

In conclusion, the evaluation of adversarial robustness is a critical area of research, requiring the development of effective metrics and evaluation frameworks. The development of effective metrics for evaluating adversarial robustness will remain a critical area of research, enabling the evaluation and improvement of defense mechanisms against adversarial attacks \cite{ref145}.

The development of effective metrics for evaluating adversarial robustness is an ongoing area of research, with many new metrics and evaluation frameworks being proposed in recent years \cite{ref181}. The importance of evaluating adversarial robustness cannot be overstated, and the development of effective metrics and evaluation frameworks is essential for advancing the field of adversarial robustness \cite{ref182}.

In summary, the evaluation of adversarial robustness is a complex and multifaceted problem, requiring the development of effective metrics and evaluation frameworks. The metrics mentioned above provide a comprehensive evaluation of a model's robustness, and the development of effective metrics will remain a critical area of research, enabling the evaluation and improvement of defense mechanisms against adversarial attacks \cite{ref140}.

The importance of evaluating adversarial robustness cannot be overstated, and the development of effective metrics and evaluation frameworks is essential for advancing the field of adversarial robustness \cite{ref183}. The use of effective metrics and evaluation frameworks will enable researchers to develop more robust models, identify areas for improvement, and advance the field of adversarial robustness \cite{ref184}.

In conclusion, the evaluation of adversarial robustness is a critical area of research, requiring the development of effective metrics and evaluation frameworks. The development of effective metrics for evaluating adversarial robustness will remain a critical area of research, enabling the evaluation and improvement of defense mechanisms against adversarial attacks \cite{ref185}.

The development of effective metrics for evaluating adversarial robustness is an ongoing area of research, with many new metrics and evaluation frameworks being proposed in recent years \cite{ref186}. The importance of evaluating adversarial robustness cannot be overstated, and the development of effective metrics and evaluation frameworks is essential for advancing the field of adversarial robustness \cite{ref187}.

In summary, the evaluation of adversarial robustness is a complex and multifaceted problem, requiring the development of effective metrics and evaluation frameworks. The metrics mentioned above provide a comprehensive evaluation of a model's robustness, and the development of effective metrics will remain a critical area of research, enabling the evaluation and improvement of defense mechanisms against adversarial attacks \cite{ref188}.

The importance of evaluating adversarial robustness cannot be overstated, and the development of effective metrics and evaluation frameworks is essential for advancing the field of adversarial robustness \cite{ref189}. The use of effective metrics and evaluation frameworks will enable researchers to develop more robust models, identify areas for improvement, and advance the field of adversarial robustness \cite{ref190}.

In conclusion, the evaluation of adversarial robustness is a critical area of research, requiring the development of effective metrics and evaluation frameworks. The development of effective metrics for evaluating adversarial robustness will remain a critical area of research, enabling the evaluation and improvement of defense mechanisms against adversarial attacks \cite{ref144}.

The development of effective metrics for evaluating adversarial robustness is an ongoing area of research, with many new metrics and evaluation frameworks being proposed in recent years \cite{ref177}. The importance of evaluating adversarial robustness cannot be overstated, and the development of effective metrics and evaluation frameworks is essential for advancing the field of adversarial robustness \cite{ref178}.

In summary, the evaluation of adversarial robustness is a complex and multifaceted problem, requiring the development of effective metrics and evaluation frameworks. The metrics mentioned above provide a comprehensive evaluation of a model's robustness, and the development of effective metrics will remain a critical area of research, enabling the evaluation and improvement of defense mechanisms against adversarial attacks \cite{ref179}.

The importance of evaluating adversarial robustness cannot be overstated, and the development of effective metrics and evaluation frameworks is essential for advancing the field of adversarial robustness \cite{ref180}. The use of effective metrics and evaluation frameworks will enable researchers to develop more robust models, identify areas for improvement, and advance the field of adversarial robustness \cite{ref145}.

In conclusion, the evaluation of adversarial robustness is a critical area of research, requiring the development of effective metrics and evaluation frameworks. The development of effective metrics for evaluating adversarial robustness will remain a critical area of research, enabling the evaluation and improvement of defense mechanisms against adversarial attacks \cite{ref181}.

The development of effective metrics for evaluating adversarial robustness is an ongoing area of research, with many new metrics and evaluation frameworks being proposed in recent years \cite{ref182}. The importance of evaluating adversarial robustness cannot be overstated, and the development of effective metrics and evaluation frameworks is essential for advancing the field of adversarial robustness \cite{ref140}.

In summary, the evaluation of adversarial robustness is a complex and multifaceted problem, requiring the development of effective metrics and evaluation frameworks. The metrics mentioned above provide a comprehensive evaluation of a model's robustness, and the development of effective metrics will remain a critical area of research, enabling the evaluation and improvement of defense mechanisms against adversarial attacks \cite{ref183}.

The importance of evaluating adversarial robustness cannot be overstated, and the development of effective metrics and evaluation frameworks is essential for advancing the field of adversarial robustness \cite{ref184}. The use of effective metrics and evaluation frameworks will enable researchers to develop more robust models, identify areas for improvement, and advance the field of adversarial robustness \cite{ref185}.

In conclusion, the evaluation of adversarial robustness is a critical area of research, requiring the development of effective metrics and evaluation frameworks. The development of effective metrics for evaluating adversarial robustness will remain a critical area of research, enabling the evaluation and improvement of defense mechanisms against adversarial attacks \cite{ref186}.

The development of effective metrics for evaluating adversarial robustness is an ongoing area of research, with many new metrics and evaluation frameworks being proposed in recent years \cite{ref187}. The importance of evaluating adversarial robustness cannot be overstated, and the development of effective metrics and evaluation frameworks is essential for advancing the field of adversarial robustness \cite{ref188}.

In summary, the evaluation of adversarial robustness is a complex and multifaceted problem, requiring the development of effective metrics and evaluation frameworks. The metrics mentioned above provide a comprehensive evaluation of a model's robustness, and the development of effective metrics will remain a critical area of research, enabling the evaluation and improvement of defense mechanisms against adversarial attacks \cite{ref189}.

The importance of evaluating adversarial robustness cannot be overstated, and the development of effective metrics and evaluation frameworks is essential for advancing the field of adversarial robustness \cite{ref190}. The use of effective metrics and evaluation frameworks will enable researchers to develop more robust models, identify areas for improvement, and advance the field of adversarial robustness \cite{ref144}.

In conclusion, the evaluation of adversarial robustness is a critical area of research, requiring the development of effective metrics and evaluation frameworks. The development of effective metrics for evaluating adversarial robustness will remain a critical area of research, enabling the evaluation and improvement of defense mechanisms against adversarial attacks \cite{ref177}.

The development of effective metrics for evaluating adversarial robustness is an ongoing area of research, with many new metrics and evaluation frameworks being proposed in recent years \cite{ref178}. The importance of evaluating adversarial robustness cannot be overstated, and the development of effective metrics and evaluation frameworks is essential for advancing the field of adversarial robustness \cite{ref179}.

In summary, the evaluation of adversarial robustness is a complex and multifaceted problem, requiring the development of effective metrics and evaluation frameworks. The metrics mentioned above provide a comprehensive evaluation of a model's robustness, and the development of effective metrics will remain a critical area of research, enabling the evaluation and improvement of defense mechanisms against adversarial attacks \cite{ref180}.

The importance of evaluating adversarial robustness cannot be overstated, and the development of effective metrics and evaluation frameworks is essential for advancing the field of adversarial robustness \cite{ref145}. The use of effective metrics and evaluation frameworks will enable researchers to develop more robust models, identify areas for improvement, and advance the field of adversarial robustness \cite{ref181}.

In conclusion, the evaluation of adversarial robustness is a critical area of research, requiring the development of effective metrics and evaluation frameworks. The development of effective metrics for evaluating adversarial robustness will remain a critical area of research, enabling the evaluation and improvement of defense mechanisms against adversarial attacks \cite{ref182}.

The development of effective metrics for evaluating adversarial robustness is an ongoing area of research, with many new metrics and evaluation frameworks being proposed in recent years \cite{ref140}. The importance of evaluating adversarial robustness cannot be overstated, and the development of effective metrics and evaluation frameworks is essential for advancing the field of adversarial robustness \cite{ref183}.

In summary, the evaluation of adversarial robustness is a complex and multifaceted problem, requiring the development of effective metrics and evaluation frameworks. The metrics mentioned above provide a comprehensive evaluation of a model's robustness, and the development of effective metrics will remain a critical area of research, enabling the evaluation and improvement of defense mechanisms against adversarial attacks \cite{ref184}.

The importance of evaluating adversarial robustness cannot be overstated, and the development of effective metrics and evaluation frameworks is essential for advancing the field of adversarial robustness \cite{ref185}. The use of effective metrics and evaluation frameworks will enable researchers to develop more robust models, identify areas for improvement, and advance the field of adversarial robustness \cite{ref186}.

In conclusion, the evaluation of adversarial robustness is a critical area of research, requiring the development of effective metrics and evaluation frameworks. The development of effective metrics for evaluating adversarial robustness will remain a critical area of research, enabling the evaluation and improvement of defense mechanisms against adversarial attacks \cite{ref187}.

The development of effective metrics for evaluating adversarial robustness is an ongoing area of research, with many new metrics and evaluation frameworks being proposed in recent years \cite{ref188}. The importance of evaluating adversarial robustness cannot be overstated, and the development of effective metrics and evaluation frameworks is essential for advancing the field of adversarial robustness \cite{ref189}.

In summary, the evaluation of adversarial robustness is a complex and multifaceted problem, requiring the development of effective metrics and evaluation frameworks. The metrics mentioned above provide a comprehensive evaluation of a model's robustness, and the development of effective metrics will remain a critical area of research, enabling the evaluation and improvement of defense mechanisms against adversarial attacks \cite{ref190}.

The importance of evaluating adversarial robustness cannot be overstated, and the development of effective metrics and evaluation frameworks is essential for advancing the field of adversarial robustness \cite{ref144}. The use of effective metrics and evaluation frameworks will enable researchers to develop more robust models, identify areas for improvement, and advance the field of adversarial robustness \cite{ref177}.

In conclusion, the evaluation of adversarial robustness is a critical area of research, requiring the development of effective metrics and evaluation frameworks. The development of effective metrics for evaluating adversarial robustness will remain a critical area of research, enabling the evaluation and improvement of defense mechanisms against adversarial attacks \cite{ref178}.

The development of effective metrics for evaluating adversarial robustness is an ongoing area of research, with many new metrics and evaluation frameworks being proposed in recent years \cite{ref179}. The importance of evaluating adversarial robustness cannot be overstated, and the development of effective metrics and evaluation frameworks is essential for advancing the field of adversarial robustness \cite{ref180}.

In summary, the evaluation of adversarial robustness is a complex and multifaceted problem, requiring the development of effective metrics and evaluation frameworks. The metrics mentioned above provide a comprehensive evaluation of a model's robustness, and the development of effective metrics will remain a critical area of research, enabling the evaluation and improvement of defense mechanisms against adversarial attacks \cite{ref145}.

The importance of evaluating adversarial robustness cannot be overstated, and the development of effective metrics and evaluation frameworks is essential for advancing the field of adversarial robustness \cite{ref181}. The use of effective metrics and evaluation frameworks will enable researchers to develop more robust models, identify areas for improvement, and advance the field of adversarial robustness \cite{ref182}.

In conclusion, the evaluation of adversarial robustness is a critical area of research, requiring the development of effective metrics and evaluation frameworks. The development of effective metrics for evaluating adversarial robustness will remain a critical area of research, enabling the evaluation and improvement of defense mechanisms against adversarial attacks \cite{ref140}.

The development of effective metrics for evaluating adversarial robustness is an ongoing area of research, with many new metrics and evaluation frameworks being proposed in recent years \cite{ref183}. The importance of evaluating adversarial robustness cannot be overstated, and the development of effective metrics and evaluation frameworks is essential for advancing the field of adversarial robustness \cite{ref184}.

In summary, the evaluation of adversarial robustness is a complex and multifaceted problem, requiring the development of effective metrics and evaluation frameworks. The metrics mentioned above provide a comprehensive evaluation of a model's robustness, and the development of effective metrics will remain a critical area of research, enabling the evaluation and improvement of defense mechanisms against adversarial attacks \cite{ref185}.

The importance of evaluating adversarial robustness cannot be overstated, and the development of effective metrics and evaluation frameworks is essential for advancing the field of adversarial robustness \cite{ref186}. The use of effective metrics and evaluation frameworks will enable researchers to develop more robust models, identify areas for improvement, and advance the field of adversarial robustness \cite{ref187}.

In conclusion, the evaluation of adversarial robustness is a critical area of research, requiring the development of effective metrics and evaluation frameworks. The development of effective metrics for evaluating adversarial robustness will remain a critical area of research, enabling the evaluation and improvement of defense mechanisms against adversarial attacks \cite{ref188}.

The development of effective metrics for evaluating adversarial robustness is an ongoing area of research, with many new metrics and evaluation frameworks being proposed in recent years \cite{ref189}. The importance of evaluating adversarial robustness cannot be overstated, and the development of effective metrics and evaluation frameworks is essential for advancing the field of adversarial robustness \cite{ref190}.

In summary, the evaluation of adversarial robustness is a complex and multifaceted problem, requiring the development of effective metrics and evaluation frameworks. The metrics mentioned above provide a comprehensive evaluation of a model's robustness, and the development of effective metrics will remain a critical area of research, enabling the evaluation and improvement of defense mechanisms against adversarial attacks \cite{ref144}.

The importance of evaluating adversarial robustness cannot be overstated, and the development of effective metrics and evaluation frameworks is essential for advancing the field of adversarial robustness \cite{ref177}. The use of effective metrics and evaluation frameworks will enable researchers to develop more robust models, identify areas for improvement, and advance the field of adversarial robustness \cite{ref178}.

In conclusion, the evaluation of adversarial robustness is a critical area of research, requiring the development of effective metrics and evaluation frameworks. The development of effective metrics for evaluating adversarial robustness will remain a critical area of research, enabling the evaluation and improvement of defense mechanisms against adversarial attacks \cite{ref179}.

The development of effective metrics for evaluating adversarial robustness is an ongoing area of research, with many new metrics and evaluation frameworks being proposed in recent years \cite{ref180}. The importance of evaluating adversarial robustness cannot be overstated, and the development of effective metrics and evaluation frameworks is essential for advancing the field of adversarial robustness \cite{ref145}.

In summary, the evaluation of adversarial robustness is a complex and multifaceted problem, requiring the development of effective metrics and evaluation frameworks. The metrics mentioned above provide a comprehensive evaluation of a model's robustness, and the development of effective metrics will remain a critical area of research, enabling the evaluation and improvement of defense mechanisms against adversarial attacks \cite{ref181}.

The importance of evaluating adversarial robustness cannot be overstated, and the development of effective metrics and evaluation frameworks is essential for advancing the field of adversarial robustness \cite{ref182}. The use of effective metrics and evaluation frameworks will enable researchers to develop more robust models, identify areas for improvement, and advance the field of adversarial robustness \cite{ref140}.

In conclusion, the evaluation of adversarial robustness is a critical area of research, requiring the development of effective metrics and evaluation frameworks. The development of effective metrics for evaluating adversarial robustness will remain a critical area of research, enabling the evaluation and improvement of defense mechanisms against adversarial attacks \cite{ref183}.

The development of effective metrics for evaluating adversarial robustness is an ongoing area of research, with many new metrics and evaluation frameworks being proposed in recent years \cite{ref184}. The importance of evaluating adversarial robustness cannot be overstated, and the development of effective metrics and evaluation frameworks is essential for advancing the field of adversarial robustness \cite{ref185}.

In summary, the evaluation of adversarial robustness is a complex and multifaceted problem, requiring the development of effective metrics and evaluation frameworks. The metrics mentioned above provide a comprehensive evaluation of a model's robustness, and the development of effective metrics will remain a critical area of research, enabling the evaluation and improvement of defense mechanisms against adversarial attacks \cite{ref186}.

The importance of evaluating adversarial robustness cannot be overstated, and the development of effective metrics and evaluation frameworks is essential for advancing the field of adversarial robustness \cite{ref187}. The use of effective metrics and evaluation frameworks will enable researchers to develop more robust models, identify areas for improvement, and advance the field of adversarial robustness \cite{ref188}.

In conclusion, the evaluation of adversarial robustness is a critical area of research, requiring the development of effective metrics and evaluation frameworks. The development of effective metrics for evaluating adversarial robustness will remain a critical area of research, enabling the evaluation and improvement of defense mechanisms against adversarial attacks \cite{ref189}.

The development of effective metrics for evaluating adversarial robustness is an ongoing area of research, with many new metrics and evaluation frameworks being proposed in recent years \cite{ref190}. The importance of evaluating adversarial robustness cannot be overstated, and the development of effective metrics and evaluation frameworks is essential for advancing the field of adversarial robustness \cite{ref144}.

In summary, the evaluation of adversarial robustness is a complex and multifaceted problem, requiring the development of effective metrics and evaluation frameworks. The metrics mentioned above provide a comprehensive evaluation of a model's robustness, and the development of effective metrics will remain a critical area of research, enabling the evaluation and improvement of defense mechanisms against adversarial attacks \cite{ref177}.

The importance of evaluating adversarial robustness cannot be overstated, and the development of effective metrics and evaluation frameworks is essential for advancing the field of adversarial robustness \cite{ref178}. The use of effective metrics and evaluation frameworks will enable researchers to develop more robust models, identify areas for improvement, and advance the field of adversarial robustness \cite{ref179}.

In conclusion, the evaluation of adversarial robustness is a critical area of research, requiring the development of effective metrics and evaluation frameworks. The development of effective metrics for evaluating adversarial robustness will remain a critical area of research, enabling the evaluation and improvement of defense mechanisms against adversarial attacks \cite{ref180}.

The development of effective metrics for evaluating adversarial robustness is an ongoing area of research, with many new metrics and evaluation frameworks being proposed in recent years \cite{ref145}. The importance of evaluating adversarial robustness cannot be overstated, and the development of effective metrics and evaluation frameworks is essential for advancing the field of adversarial robustness \cite{ref181}.

In summary, the evaluation of adversarial robustness is a complex and multifaceted problem, requiring the development of effective metrics and evaluation frameworks. The metrics mentioned above provide a comprehensive evaluation of a model's robustness, and the development of effective metrics will remain a critical area of research, enabling the evaluation and improvement of defense mechanisms against adversarial attacks \cite{ref182}.

The importance of evaluating adversarial robustness cannot be overstated, and the development of effective metrics and evaluation frameworks is essential for advancing the field of adversarial robustness \cite{ref140}. The use of effective metrics and evaluation frameworks will enable researchers to develop more robust models, identify areas for improvement, and advance the field of adversarial robustness \cite{ref183}.

In conclusion, the evaluation of adversarial robustness is a critical area of research, requiring the development of effective metrics and evaluation frameworks. The development of effective metrics for evaluating adversarial robustness will remain a critical area of research, enabling the evaluation and improvement of defense mechanisms against adversarial attacks \cite{ref184}.

The development of effective metrics for evaluating adversarial robustness is an ongoing area of research, with many new metrics and evaluation frameworks being proposed in recent years \cite{ref185}. The importance of evaluating adversarial robustness cannot be overstated, and the development of effective metrics and evaluation frameworks is essential for advancing the field of adversarial robustness \cite{ref186}.

In summary, the evaluation of adversarial robustness is a complex and multifaceted problem, requiring the development of effective metrics and evaluation frameworks. The metrics mentioned above provide a comprehensive evaluation of a model's robustness, and the development of effective metrics will remain a critical area of research, enabling the evaluation and improvement of defense mechanisms against adversarial attacks \cite{ref187}.

The importance of evaluating adversarial robustness cannot be overstated, and the development of effective metrics and evaluation frameworks is essential for advancing the field of adversarial robustness \cite{ref188}. The use of effective metrics and evaluation frameworks will enable researchers to develop more robust models, identify areas for improvement, and advance the field of adversarial robustness \cite{ref189}.

In conclusion, the evaluation of adversarial robustness is a critical area of research, requiring the development of effective metrics and evaluation frameworks. The development of effective metrics for evaluating adversarial robustness will remain a critical area of research, enabling the evaluation and improvement of defense mechanisms against adversarial attacks \cite{ref190}.

The development of effective metrics for evaluating adversarial robustness is an ongoing area of research, with many new metrics and evaluation frameworks being proposed in recent years \cite{ref144}. The importance of evaluating adversarial robustness cannot be overstated, and the development of effective metrics and evaluation frameworks is essential for advancing the field of adversarial robustness \cite{ref177}.

In summary, the evaluation of adversarial robustness is a complex and multifaceted problem, requiring the development of effective metrics and evaluation frameworks. The metrics mentioned above provide a comprehensive evaluation of a model's robustness, and the development of effective metrics will remain a critical area of research, enabling the evaluation and improvement of defense mechanisms against adversarial attacks \cite{ref178}.

The importance of evaluating adversarial robustness cannot be overstated, and the development of effective metrics and evaluation frameworks is essential for advancing the field of adversarial robustness \cite{ref179}. The use of effective metrics and evaluation frameworks will enable researchers

\subsection{Benchmarking and Evaluation Frameworks}

Benchmarking and evaluation frameworks play a crucial role in assessing the performance of computer vision models under adversarial attacks. These frameworks provide a standardized way to evaluate the robustness of models, enabling researchers to compare and contrast different defense mechanisms. In this subsection, we will discuss existing benchmarking and evaluation frameworks for adversarial robustness, including datasets, protocols, and tools used to assess the performance of computer vision models.

One of the most widely used benchmarking frameworks for adversarial robustness is RobustBench \cite{ref191}. RobustBench provides a comprehensive evaluation of the adversarial robustness of image classification models on the CIFAR-10 dataset. The framework uses AutoAttack, an ensemble of white-box and black-box attacks, to evaluate the robustness of models. RobustBench also provides a leaderboard that ranks models based on their robustness, making it easier for researchers to compare and identify the most robust models.

Other notable benchmarking frameworks include ImageNet-Patch \cite{ref192}, which focuses on evaluating the robustness of models against adversarial patches, and FACESEC \cite{ref166}, which provides a fine-grained evaluation of the robustness of face recognition systems. Additionally, NIC-RobustBench \cite{ref193} provides a comprehensive evaluation of the robustness of neural image compression models.

The use of benchmarking and evaluation frameworks is essential for developing more robust computer vision models. By providing a standardized way to evaluate the robustness of models, these frameworks can help researchers identify the most effective defense mechanisms, develop more robust models, and advance the field of computer vision. For example, \cite{ref194} shows that the use of meticulous adversarial attacks can be effective in evaluating the robustness of vision-language models.

In addition to these frameworks, researchers have also proposed new frameworks and datasets for evaluating the robustness of computer vision models. For instance, \cite{ref152} proposes a new framework for evaluating the robustness of vision-based reinforcement learning models, while \cite{ref195} proposes a new framework for evaluating the robustness of machine learning models against adversarial attacks.

The development of benchmarking and evaluation frameworks is an active area of research, with new frameworks and datasets being proposed regularly. As the field of computer vision continues to evolve, the development of benchmarking and evaluation frameworks will play a critical role in ensuring the security and reliability of computer vision systems in real-world applications. As shown in \cite{ref162}, the development of robust models is critical for ensuring the security and reliability of computer vision systems in real-world applications.

In conclusion, benchmarking and evaluation frameworks are essential for developing more robust computer vision models. By providing a standardized way to evaluate the robustness of models, these frameworks can help researchers identify the most effective defense mechanisms, develop more robust models, and advance the field of computer vision. The use of benchmarking and evaluation frameworks can also help researchers develop more robust models by identifying the weaknesses of existing models and providing a set of protocols and tools for evaluating the robustness of models. As the field of computer vision continues to evolve, the development of benchmarking and evaluation frameworks will play a critical role in ensuring the security and reliability of computer vision systems in real-world applications \cite{ref175}.

\section{Applications and Case Studies of Physical-World Adversarial Attacks and Defenses}

\subsection{Autonomous Vehicle Systems}

Autonomous vehicle systems have become a crucial area of research in the field of computer vision, with the primary goal of enabling vehicles to operate safely and efficiently without human intervention. However, these systems are vulnerable to visual adversarial attacks, which can compromise their performance and pose significant risks to road safety \cite{ref196}. In this subsection, we will discuss the applications of visual adversarial attacks and defenses in autonomous vehicle systems, highlighting the key challenges and potential solutions.

Visual adversarial attacks on autonomous vehicle systems can take various forms, including attacks on object detection, semantic segmentation, and image classification \cite{ref63}. For instance, \cite{ref197} demonstrates the vulnerability of deep neural networks to adversarial attacks in object detection, which is a critical component of autonomous vehicle systems. Another example is the \cite{ref198} paper, which demonstrates the vulnerability of deep learning-based automated lane centering systems to physical-world adversarial attacks \cite{ref198}.

To defend against these visual adversarial attacks, various defense mechanisms have been proposed, including adversarial training, input preprocessing, and detection methods \cite{ref197}. For example, \cite{ref199} proposes a robust deep reinforcement learning approach to enhance the security and safety of autonomous vehicle systems. Additionally, the use of multiple sensors and sensor fusion models can provide redundant and complementary information to enhance the security and safety of autonomous vehicle systems \cite{ref200}.

Several papers have investigated the vulnerability of autonomous vehicle systems to visual adversarial attacks, including \cite{ref201}, \cite{ref202}, and \cite{ref203}. These papers demonstrate the effectiveness of various attack methods in compromising the performance of autonomous vehicle systems. Furthermore, papers such as \cite{ref204} and \cite{ref33} have proposed defense mechanisms to mitigate the impact of visual adversarial attacks on autonomous vehicle systems.

Other notable papers in this area include \cite{ref205}, \cite{ref206}, and \cite{ref207}. These papers provide a comprehensive overview of the current state of research in visual adversarial attacks and defenses in autonomous vehicle systems, highlighting the need for continued innovation and development in this critical area.

In conclusion, visual adversarial attacks pose a significant threat to the security and safety of autonomous vehicle systems \cite{ref196}. Various types of attacks, including attacks on object detection, semantic segmentation, and lane tracking, can compromise the performance of these systems and pose significant risks to road safety. To defend against these attacks, various defense mechanisms have been proposed, including adversarial training, input preprocessing, and detection methods. Further research is needed to develop more effective defense mechanisms and to enhance the security and safety of autonomous vehicle systems against visual adversarial attacks \cite{ref63}.

\subsection{Surveillance Systems}

Surveillance systems have become an essential component of modern security infrastructure, leveraging computer vision and deep learning techniques to detect, track, and recognize objects and individuals. As we have seen in the previous subsection, autonomous vehicle systems are vulnerable to visual adversarial attacks, and similarly, surveillance systems are also susceptible to such attacks, posing a significant threat to their effectiveness and reliability. In this subsection, we explore the applications of visual adversarial attacks and defenses in surveillance systems, highlighting the key challenges and potential solutions.

Recent studies have demonstrated the feasibility of launching physical adversarial attacks on surveillance systems, including \cite{ref11} and \cite{ref63}. These attacks can be used to evade detection, tracking, or recognition by manipulating the input data in a way that is imperceptible to humans. For instance, an attacker can wear a specially designed t-shirt or glasses to avoid being detected by a surveillance camera \cite{ref32}. Similarly, \cite{ref29} demonstrates how malicious ads or logos can be used to deceive traffic sign recognition systems.

To defend against such attacks, various techniques have been proposed, including adversarial training, input preprocessing, and detection methods \cite{ref208}. Adversarial training involves training the model on adversarial examples to improve its robustness, while input preprocessing techniques aim to detect and remove adversarial perturbations from the input data. Detection methods, on the other hand, focus on identifying and flagging potential adversarial attacks. For example, \cite{ref209} proposes a deep learning-based approach to enhance the security of real-time video surveillance systems against adversarial attacks.

Furthermore, \cite{ref210} presents a real-time robust video object detection system that can defend against physical-world adversarial attacks. Additionally, \cite{ref211} investigates the vulnerability of video anomaly detection systems to adversarial machine learning attacks, and \cite{ref33} proposes a moving target defense approach to mitigate such attacks.

Other notable studies in this area include \cite{ref212}, which presents a black-box adversarial attack scheme based on semantic segmentation and model inversion (SSMI) that can make objects of interest disappear without accessing object detectors. Moreover, \cite{ref213} introduces a distributed video surveillance system that leverages artificial intelligence and digital twins technologies to enhance the security and efficiency of surveillance systems.

In conclusion, the applications of visual adversarial attacks and defenses in surveillance systems are a critical area of research, with significant implications for the security and reliability of these systems. As we will see in the following subsection, the vulnerability of facial recognition systems to visual adversarial attacks is also a significant concern, and ongoing efforts are needed to develop more effective defense strategies and to improve the robustness of these systems to adversarial attacks \cite{ref11}. By understanding the vulnerabilities of surveillance systems to adversarial attacks and developing effective defense mechanisms, we can enhance the overall security and efficiency of these systems, and ultimately improve the safety and reliability of various applications that rely on computer vision and deep learning techniques.

\subsection{Facial Recognition Systems}

Facial recognition systems have become increasingly prevalent in various aspects of our lives, from security and surveillance to social media and mobile devices. As we have seen in the previous subsection, surveillance systems are vulnerable to visual adversarial attacks, and facial recognition systems are no exception. In fact, the susceptibility of facial recognition systems to visual adversarial attacks is a significant concern, as these systems are often used in security-critical applications. In this subsection, we will examine the applications of visual adversarial attacks and defenses in facial recognition systems, and explore the ways in which these attacks can be used to compromise the security and reliability of these systems.

Visual adversarial attacks on facial recognition systems can take many forms, including digital and physical attacks. Digital attacks involve manipulating the input image to create an adversarial example that can fool the facial recognition system \cite{ref214}. These attacks can be launched using various techniques, such as adding noise to the image or using generative adversarial networks (GANs) to create synthetic faces \cite{ref215}. Physical attacks, on the other hand, involve manipulating the physical environment to create an adversarial example that can fool the facial recognition system \cite{ref36}. For example, an attacker could use a physical adversarial attack to create a mask or disguise that can fool a facial recognition system into identifying an individual as someone else \cite{ref216}.

The potential consequences of visual adversarial attacks on facial recognition systems are severe, and include identity theft, spoofing, and other forms of malicious activity \cite{ref217}. To defend against these types of attacks, various techniques have been proposed, including adversarial training, input preprocessing, and detection-based methods \cite{ref6}. Adversarial training involves training the facial recognition system on a dataset that includes adversarial examples, which can help improve its robustness to attacks \cite{ref218}. Input preprocessing techniques, such as image filtering or normalization, can also be used to reduce the effectiveness of adversarial attacks \cite{ref128}. Detection-based methods, on the other hand, involve detecting and responding to adversarial attacks in real-time, using techniques such as anomaly detection or machine learning-based classification \cite{ref219}.

In addition to these techniques, various other methods have been proposed to defend against visual adversarial attacks on facial recognition systems. For example, \cite{ref220} proposes a network that can alleviate adversarial perturbations in facial images, while \cite{ref166} proposes a framework for evaluating the robustness of facial recognition systems to adversarial attacks. These methods can be used in conjunction with other defense strategies to provide a robust and comprehensive defense against visual adversarial attacks.

As we will see in the following subsection, the applications of visual adversarial attacks and defenses are not limited to facial recognition systems, but can also be applied to other domains, such as autonomous vehicles and healthcare. However, the unique characteristics of facial recognition systems make them a critical area of research, and ongoing efforts are needed to develop more effective defense strategies and to improve the robustness of these systems to adversarial attacks \cite{ref11}. By examining the applications of visual adversarial attacks and defenses in facial recognition systems, we can better understand the risks and challenges associated with these systems and develop more effective strategies for mitigating these risks \cite{ref215}.

\section{Future Directions, Open Challenges, and Conclusion}

\subsection{Introduction to Future Directions}

The field of visual adversarial attacks and defenses has witnessed significant advancements in recent years, with a growing body of research focused on understanding and mitigating the vulnerabilities of deep learning models to adversarial perturbations. As we look to the future, it is essential to identify the key directions that will shape the development of this field. In this subsection, we introduce the concept of future directions in visual adversarial attacks and defenses, highlighting the emerging trends, challenges, and opportunities that will drive innovation in this area.

One of the primary future directions in visual adversarial attacks is the development of more sophisticated and realistic attack methods. Current attack methods, such as \cite{ref104} and \cite{ref101}, have demonstrated impressive results in fooling deep learning models. However, these attacks are often limited to specific scenarios or datasets, and there is a need for more generalizable and robust attack methods that can be applied to a wide range of applications. Future research should focus on developing attack methods that can adapt to different environments, lighting conditions, and other real-world factors.

Another critical direction is the development of more effective defense mechanisms against visual adversarial attacks. While significant progress has been made in this area, with techniques such as \cite{ref111} and \cite{ref221} showing promise, there is still a need for more robust and generalizable defense methods. Future research should explore new architectures, training methods, and regularization techniques that can improve the robustness of deep learning models to adversarial attacks. Additionally, the development of more effective evaluation metrics and benchmarking frameworks, such as \cite{ref90}, will be crucial in assessing the effectiveness of defense mechanisms and identifying areas for improvement.

The emergence of new technologies, such as \cite{ref222} and \cite{ref152}, is also expected to play a significant role in shaping the future of visual adversarial attacks and defenses. These technologies have the potential to enable more sophisticated and realistic attack methods, as well as more effective defense mechanisms. For example, \cite{ref222} has demonstrated the ability to generate highly realistic and effective adversarial attacks on 3D models, while \cite{ref223} has shown promise in improving the robustness of vision-based reinforcement learning models to adversarial attacks.

Furthermore, the increasing use of deep learning models in real-world applications, such as autonomous vehicles, surveillance systems, and facial recognition systems, highlights the need for more practical and effective defense mechanisms. Future research should focus on developing defense mechanisms that can be deployed in real-world scenarios, taking into account the constraints and challenges of these applications. This may involve developing more efficient and scalable defense methods, as well as exploring new techniques for detecting and mitigating adversarial attacks in real-time.

In addition to these technical challenges, there are also ethical and societal implications of visual adversarial attacks and defenses that need to be considered. As deep learning models become increasingly ubiquitous in our daily lives, it is essential to ensure that they are designed and developed with robustness and security in mind. This may involve developing new guidelines and regulations for the development and deployment of deep learning models, as well as raising awareness about the potential risks and vulnerabilities of these models.

Finally, the future of visual adversarial attacks and defenses will also be shaped by the development of new datasets and evaluation frameworks. The availability of high-quality datasets and evaluation frameworks, such as \cite{ref90}, will be crucial in driving innovation in this area. Future research should focus on developing new datasets and evaluation frameworks that can simulate real-world scenarios and provide a comprehensive assessment of the effectiveness of defense mechanisms.

In conclusion, the future of visual adversarial attacks and defenses is exciting and challenging, with many emerging trends, challenges, and opportunities that will drive innovation in this area. As we look to the future, it is essential to identify the key directions that will shape the development of this field, including the development of more sophisticated and realistic attack methods, more effective defense mechanisms, and more practical and effective defense mechanisms for real-world applications. By addressing these challenges and opportunities, we can ensure that deep learning models are designed and developed with robustness and security in mind, and that they can be trusted to operate effectively and safely in a wide range of applications, ultimately paving the way for the development of more secure and reliable systems that can withstand the open challenges discussed in the subsequent section.

\subsection{Open Challenges in Visual Adversarial Attacks and Defenses}

Despite the significant progress made in visual adversarial attacks and defenses, as discussed in the previous section on future directions, there are still several open challenges that need to be addressed in order to develop more robust and secure computer vision systems. One of the major challenges is the need for more sophisticated attack and defense strategies \cite{ref34}. Current attack methods, such as the Fast Gradient Sign Method (FGSM) and Projected Gradient Descent (PGD), are effective but can be detected and defended against using various defense mechanisms \cite{ref6}. However, more advanced attack methods, such as the Carlini \& Wagner (C\&W) attack, can still bypass these defenses \cite{ref224}. This highlights the importance of developing more robust and generalizable defense mechanisms that can defend against a wide range of attacks.

Another challenge is the lack of robustness of current defense mechanisms against unknown or unseen attacks \cite{ref225}. Most defense mechanisms are designed to defend against specific types of attacks, but they may not be effective against other types of attacks \cite{ref196}. Therefore, there is a need for more robust and generalizable defense mechanisms that can defend against a wide range of attacks \cite{ref109}. This can be achieved through the development of more advanced attack and defense strategies, as well as a deeper understanding of the underlying mechanisms of adversarial attacks.

The development of more sophisticated attack and defense strategies is also hindered by the lack of standardization in the field \cite{ref176}. Different studies use different metrics and evaluation protocols, making it difficult to compare and contrast the effectiveness of different attack and defense methods \cite{ref145}. Therefore, there is a need for standardization in the field to facilitate the development of more effective attack and defense strategies \cite{ref226}. This standardization can help to ensure that defense mechanisms are evaluated and compared in a consistent and fair manner, which is essential for developing more robust and secure computer vision systems.

Furthermore, the increasing use of deep learning models in real-world applications, such as autonomous vehicles and surveillance systems, has raised concerns about the security and reliability of these systems \cite{ref63}. Adversarial attacks can have significant consequences in these applications, and therefore, there is a need for more robust and secure defense mechanisms \cite{ref8}. This highlights the importance of addressing the open challenges in visual adversarial attacks and defenses, which will be crucial in ensuring the safety and security of computer vision systems in real-world applications.

In addition, the development of more sophisticated attack and defense strategies is also limited by the lack of understanding of the underlying mechanisms of adversarial attacks \cite{ref227}. Current attack methods are often based on heuristics and empirical observations, rather than a deep understanding of the underlying mechanisms \cite{ref105}. Therefore, there is a need for more research on the underlying mechanisms of adversarial attacks to develop more effective and robust defense mechanisms \cite{ref228}. This research can help to shed light on the underlying mechanisms of adversarial attacks, which can inform the development of more robust and secure defense mechanisms.

The use of multimodal data, such as images and text, has also raised new challenges in visual adversarial attacks and defenses \cite{ref229}. Current attack methods are often designed for unimodal data, and therefore, there is a need for more research on multimodal attack and defense methods \cite{ref230}. Furthermore, the increasing use of transfer learning and few-shot learning has also raised new challenges in visual adversarial attacks and defenses \cite{ref231}. These challenges highlight the need for more research on the development of robust and secure computer vision systems that can operate effectively in a wide range of scenarios.

Finally, the development of more sophisticated attack and defense strategies is also limited by the lack of availability of large-scale datasets and computational resources \cite{ref232}. Current datasets and computational resources are often limited, and therefore, there is a need for more research on large-scale datasets and computational resources to develop more effective and robust defense mechanisms \cite{ref233}. By addressing these challenges and developing more robust and secure computer vision systems, we can mitigate the risks associated with visual adversarial attacks and ensure the safe and reliable operation of computer vision systems in real-world applications, as will be discussed in the following section.

\subsection{Conclusion and Recommendations}

In conclusion, this survey has provided a comprehensive overview of the current state of visual adversarial attacks and defenses in the physical world, highlighting the importance of considering these threats in the development and deployment of computer vision systems. The key findings and takeaways from this survey emphasize the need for practitioners, researchers, and policymakers to prioritize the security and reliability of these systems in real-world applications, as demonstrated by various studies, including \cite{ref63} and \cite{ref99}.

Building on the challenges and open research directions discussed in the previous section, it is clear that physical adversarial attacks can have significant consequences, such as compromising the safety and security of autonomous vehicles, surveillance systems, and facial recognition systems. The survey has also highlighted the various types of physical adversarial attacks, including physical, digital, and hybrid attacks, and their characteristics, advantages, and limitations. For instance, \cite{ref32} has shown that physical adversarial attacks can be launched using light projection, which can be used to impersonate or evade recognition. Moreover, \cite{ref29} has demonstrated that physical adversarial attacks can be used to deceive traffic sign recognition systems, which can have significant safety implications.

To address these challenges and enhance the security and reliability of computer vision systems, several recommendations can be made. Firstly, practitioners and researchers should prioritize the development of robust and secure computer vision systems that can withstand physical adversarial attacks. This can be achieved through the use of techniques such as adversarial training, input preprocessing, and gradient masking, as demonstrated by \cite{ref39} and \cite{ref234}. Secondly, policymakers should establish standards and regulations for the development and deployment of computer vision systems, including guidelines for ensuring their security and reliability.

In terms of future research directions, several areas can be identified. Firstly, there is a need for more research on the development of robust and secure computer vision systems that can withstand physical adversarial attacks. This can include the use of techniques such as adversarial training, input preprocessing, and gradient masking. Secondly, there is a need for more research on the development of effective defense strategies against physical adversarial attacks, including the use of techniques such as detection and localization of adversarial patches. Finally, there is a need for more research on the standardization and regulation of computer vision systems, including guidelines for ensuring their security and reliability, as emphasized by \cite{ref67} and \cite{ref152}.

Overall, this survey has provided a comprehensive overview of the current state of visual adversarial attacks and defenses in the physical world, highlighting the importance of considering these threats in the development and deployment of computer vision systems. By prioritizing the development of robust and secure computer vision systems, establishing standards and regulations for their development and deployment, and conducting further research in the area of physical adversarial attacks and defenses, we can enhance the security and reliability of computer vision systems and mitigate the risks associated with physical adversarial attacks. As \cite{ref34} and \cite{ref235} have emphasized, the development of robust and secure computer vision systems is crucial for ensuring the safety and security of various applications, including autonomous vehicles, surveillance systems, and facial recognition systems. By prioritizing the security and reliability of these systems, we can prevent potential attacks and ensure the safe and reliable operation of these systems.


\begin{thebibliography}{235}
\addcontentsline{toc}{section}{References}
\bibitem{ref1} State-of-the-art optical-based physical adversarial attacks for deep learning computer vision systems. arXiv:2303.12249, 2023. \url{https://arxiv.org/abs/2303.12249}
\bibitem{ref2} Detecting and Segmenting Adversarial Graphics Patterns from Images. arXiv:2108.09383, 2021. \url{https://arxiv.org/abs/2108.09383}
\bibitem{ref3} Metamorphic Detection of Adversarial Examples in Deep Learning Models with Affine Transformations. doi:10.1109/met.2019.00016, 2019. \url{https://doi.org/10.1109/met.2019.00016}
\bibitem{ref4} Threat of Adversarial Attacks on Deep Learning in Computer Vision: A Survey. doi:10.1109/access.2018.2807385, 2018. \url{https://doi.org/10.1109/access.2018.2807385}
\bibitem{ref5} Physical Adversarial Attack Meets Computer Vision: A Decade Survey. doi:10.1109/tpami.2024.3430860, 2024. \url{https://doi.org/10.1109/tpami.2024.3430860}
\bibitem{ref6} Adversarial Attacks in Computer Vision: Challenges and Defense Strategies. doi:10.63345/sjaibt.v2.i4.101, 2025. \url{https://doi.org/10.63345/sjaibt.v2.i4.101}
\bibitem{ref7} A Survey on Physical Adversarial Attack in Computer Vision. arXiv:2209.14262, 2022. \url{https://arxiv.org/abs/2209.14262}
\bibitem{ref8} Adversarial catoptric light: An effective, stealthy and robust physical-world attack to DNNs. doi:10.1049/cvi2.12264, 2024. \url{https://doi.org/10.1049/cvi2.12264}
\bibitem{ref9} Shadows can be Dangerous: Stealthy and Effective Physical-world Adversarial Attack by Natural Phenomenon. arXiv:2203.03818, 2022. \url{https://arxiv.org/abs/2203.03818}
\bibitem{ref10} Detecting Adversarial Patches with Class Conditional Reconstruction Networks. arXiv:2011.05850, 2020. \url{https://arxiv.org/abs/2011.05850}
\bibitem{ref11} Physical Adversarial Attacks for Surveillance: A Survey. doi:10.1109/tnnls.2023.3321432, 2023. \url{https://doi.org/10.1109/tnnls.2023.3321432}
\bibitem{ref12} Improving Adversarial Transferability of Visual-Language Pre-training Models through Collaborative Multimodal Interaction. arXiv:2403.10883, 2024. \url{https://arxiv.org/abs/2403.10883}
\bibitem{ref13} Invisible Perturbations: Physical Adversarial Examples Exploiting the Rolling Shutter Effect. doi:10.1109/cvpr46437.2021.01443, 2021. \url{https://doi.org/10.1109/cvpr46437.2021.01443}
\bibitem{ref14} Transferable Adversarial Attacks on Black-Box Vision-Language Models. arXiv:2505.01050, 2025. \url{https://arxiv.org/abs/2505.01050}
\bibitem{ref15} Adversarial Attacks on Deep-Learning Based Radio Signal Classification. arXiv:1808.07713, 2018. \url{https://arxiv.org/abs/1808.07713}
\bibitem{ref16} A Survey and Evaluation of Adversarial Attacks in Object Detection. doi:10.1109/tnnls.2025.3561225, 2025. \url{https://doi.org/10.1109/tnnls.2025.3561225}
\bibitem{ref17} Attack to Fool and Explain Deep Networks. arXiv:2106.10606, 2021. \url{https://arxiv.org/abs/2106.10606}
\bibitem{ref18} Physical Adversarial Examples for Object Detectors. arXiv:1807.07769, 2018. \url{https://arxiv.org/abs/1807.07769}
\bibitem{ref19} Adversarial Attacks on Monocular Pose Estimation. doi:10.1109/iros47612.2022.9982154, 2022. \url{https://doi.org/10.1109/iros47612.2022.9982154}
\bibitem{ref20} Adversarial Evasion Attacks on Computer Vision using SHAP Values. doi:10.13140/rg.2.2.28762.40647, 2024. \url{https://doi.org/10.13140/rg.2.2.28762.40647}
\bibitem{ref21} On the Real-World Adversarial Robustness of Real-Time Semantic Segmentation Models for Autonomous Driving. arXiv:2201.01850, 2022. \url{https://arxiv.org/abs/2201.01850}
\bibitem{ref22} Adversarial camera stickers: A physical camera-based attack on deep learning systems. arXiv:1904.00759, 2019. \url{https://arxiv.org/abs/1904.00759}
\bibitem{ref23} Impact of Adversarial Examples on Deep Learning Models for Biomedical Image Segmentation. doi:10.1007/978-3-030-32245-8\_34, 2019. \url{https://doi.org/10.1007/978-3-030-32245-8_34}
\bibitem{ref24} Adversarial Attacks and Defences: A Survey. arXiv:1810.00069, 2018. \url{https://arxiv.org/abs/1810.00069}
\bibitem{ref25} Adversarial Training for Deep Learning-based Intrusion Detection Systems. arXiv:2104.09852, 2021. \url{https://arxiv.org/abs/2104.09852}
\bibitem{ref26} DAPAS : Denoising Autoencoder to Prevent Adversarial attack in Semantic Segmentation. doi:10.1109/ijcnn48605.2020.9207291, 2020. \url{https://doi.org/10.1109/ijcnn48605.2020.9207291}
\bibitem{ref27} Wild Patterns: Ten Years After the Rise of Adversarial Machine Learning. doi:10.1016/j.patcog.2018.07.023, 2018. \url{https://doi.org/10.1016/j.patcog.2018.07.023}
\bibitem{ref28} Simple Physical Adversarial Examples against End-to-End Autonomous Driving Models. arXiv:1903.05157, 2019. \url{https://arxiv.org/abs/1903.05157}
\bibitem{ref29} Rogue Signs: Deceiving Traffic Sign Recognition with Malicious Ads and Logos. arXiv:1801.02780, 2018. \url{https://arxiv.org/abs/1801.02780}
\bibitem{ref30} Fooling the Eyes of Autonomous Vehicles: Robust Physical Adversarial Examples Against Traffic Sign Recognition Systems. arXiv:2201.06192, 2022. \url{https://arxiv.org/abs/2201.06192}
\bibitem{ref31} A comprehensive survey of deep face verification systems adversarial attacks and defense strategies. doi:10.1038/s41598-025-15753-8, 2025. \url{https://doi.org/10.1038/s41598-025-15753-8}
\bibitem{ref32} Adversarial Light Projection Attacks on Face Recognition Systems: A Feasibility Study. arXiv:2003.11145, 2020. \url{https://arxiv.org/abs/2003.11145}
\bibitem{ref33} Moving target defense for embedded deep visual sensing against adversarial examples. doi:10.1145/3356250.3360025, 2019. \url{https://doi.org/10.1145/3356250.3360025}
\bibitem{ref34} Adversarial Attacks and Defenses on 3D Point Cloud Classification: A Survey. doi:10.1109/access.2023.3345000, 2023. \url{https://doi.org/10.1109/access.2023.3345000}
\bibitem{ref35} Adversarial Examples for Semantic Segmentation and Object Detection. doi:10.1109/iccv.2017.153, 2017. \url{https://doi.org/10.1109/iccv.2017.153}
\bibitem{ref36} Invisible Adversarial Attacks on Deep Learning-Based Face Recognition Models. doi:10.1109/access.2023.3279488, 2023. \url{https://doi.org/10.1109/access.2023.3279488}
\bibitem{ref37} Chain of Attack: On the Robustness of Vision-Language Models Against Transfer-Based Adversarial Attacks. doi:10.1109/cvpr52734.2025.01368, 2025. \url{https://doi.org/10.1109/cvpr52734.2025.01368}
\bibitem{ref38} Recurrent Attention Model with Log-Polar Mapping is Robust against Adversarial Attacks. arXiv:2002.05388, 2020. \url{https://arxiv.org/abs/2002.05388}
\bibitem{ref39} Certified Defenses for Adversarial Patches. arXiv:2003.06693, 2020. \url{https://arxiv.org/abs/2003.06693}
\bibitem{ref40} Adversarial Example Detection by Classification for Deep Speech Recognition. doi:10.1109/icassp40776.2020.9054750, 2020. \url{https://doi.org/10.1109/icassp40776.2020.9054750}
\bibitem{ref41} Gradient-based Adversarial Attacks against Text Transformers. arXiv:2104.13733, 2021. \url{https://arxiv.org/abs/2104.13733}
\bibitem{ref42} Simple Black-Box Adversarial Examples Generation with Very Few Queries. doi:10.1587/transinf.2019inp0002, 2020. \url{https://doi.org/10.1587/transinf.2019inp0002}
\bibitem{ref43} Query-Efficient Black-box Adversarial Attacks Guided by a Transfer-based Prior. doi:10.1109/tpami.2021.3126733, 2021. \url{https://doi.org/10.1109/tpami.2021.3126733}
\bibitem{ref44} Similarity-based Gray-box Adversarial Attack Against Deep Face Recognition. doi:10.1109/fg52635.2021.9667076, 2021. \url{https://doi.org/10.1109/fg52635.2021.9667076}
\bibitem{ref45} AdvSmo: Black-box Adversarial Attack by Smoothing Linear Structure of Texture. arXiv:2206.10988, 2022. \url{https://arxiv.org/abs/2206.10988}
\bibitem{ref46} A Frank-Wolfe Framework for Efficient and Effective Adversarial Attacks. doi:10.1609/aaai.v34i04.5753, 2020. \url{https://doi.org/10.1609/aaai.v34i04.5753}
\bibitem{ref47} Adversarial Attacks for Optical Flow-Based Action Recognition Classifiers. arXiv:1811.11875, 2018. \url{https://arxiv.org/abs/1811.11875}
\bibitem{ref48} BlackboxBench: A Comprehensive Benchmark of Black-Box Adversarial Attacks. doi:10.1109/tpami.2025.3574432, 2025. \url{https://doi.org/10.1109/tpami.2025.3574432}
\bibitem{ref49} Gray-box Adversarial Training. doi:10.1007/978-3-030-01267-0\_13, 2018. \url{https://doi.org/10.1007/978-3-030-01267-0_13}
\bibitem{ref50} Improving the adversarial transferability with relational graphs ensemble adversarial attack. doi:10.3389/fnins.2022.1094795, 2023. \url{https://doi.org/10.3389/fnins.2022.1094795}
\bibitem{ref51} Defending against Adversarial Images using Basis Functions Transformations. arXiv:1803.10840, 2018. \url{https://arxiv.org/abs/1803.10840}
\bibitem{ref52} A Survey of Black-Box Adversarial Attacks on Computer Vision Models. arXiv:1912.01667, 2019. \url{https://arxiv.org/abs/1912.01667}
\bibitem{ref53} Revisiting GANs by Best-Response Constraint: Perspective, Methodology, and Application. arXiv:2205.10146, 2022. \url{https://arxiv.org/abs/2205.10146}
\bibitem{ref54} Adversarial examples for generative models. arXiv:1702.06832, 2017. \url{https://arxiv.org/abs/1702.06832}
\bibitem{ref55} Adversarial examples by perturbing high-level features in intermediate decoder layers. doi:10.5220/0010892800003116, 2022. \url{https://doi.org/10.5220/0010892800003116}
\bibitem{ref56} Efficient Optimization Algorithms for Linear Adversarial Training. arXiv:2410.12677, 2024. \url{https://arxiv.org/abs/2410.12677}
\bibitem{ref57} Generating Adversarial Examples through Latent Space Exploration of Generative Adversarial Networks. doi:10.1145/3583133.3596392, 2023. \url{https://doi.org/10.1145/3583133.3596392}
\bibitem{ref58} A Unified Framework for Adversarial Attack and Defense in Constrained Feature Space. arXiv:2112.01156, 2021. \url{https://arxiv.org/abs/2112.01156}
\bibitem{ref59} Re-purposing Heterogeneous Generative Ensembles with Evolutionary Computation. arXiv:2003.13532, 2020. \url{https://arxiv.org/abs/2003.13532}
\bibitem{ref60} They Might NOT Be Giants: Crafting Black-Box Adversarial Examples with Fewer Queries Using Particle Swarm Optimization. arXiv:1909.07490, 2019. \url{https://arxiv.org/abs/1909.07490}
\bibitem{ref61} Adversarial Machine Learning in Network Intrusion Detection Systems. doi:10.1016/j.eswa.2021.115782, 2021. \url{https://doi.org/10.1016/j.eswa.2021.115782}
\bibitem{ref62} Opportunities and Challenges in Deep Learning Adversarial Robustness: A Survey. arXiv:2007.00753, 2020. \url{https://arxiv.org/abs/2007.00753}
\bibitem{ref63} Physical Adversarial Attacks for Camera-Based Smart Systems: Current Trends, Categorization, Applications, Research Challenges, and Future Outlook. doi:10.1109/access.2023.3321118, 2023. \url{https://doi.org/10.1109/access.2023.3321118}
\bibitem{ref64} All You Need is RAW: Defending Against Adversarial Attacks with Camera Image Pipelines. doi:10.1007/978-3-031-19800-7\_19, 2022. \url{https://doi.org/10.1007/978-3-031-19800-7_19}
\bibitem{ref65} Multiclass ASMA vs Targeted PGD Attack in Image Segmentation. arXiv:2208.01844, 2022. \url{https://arxiv.org/abs/2208.01844}
\bibitem{ref66} Adversarial Deep Learning: A Survey on Adversarial Attacks and Defense Mechanisms on Image Classification. doi:10.1109/access.2022.3208131, 2022. \url{https://doi.org/10.1109/access.2022.3208131}
\bibitem{ref67} Revisiting Transferable Adversarial Images: Systemization, Evaluation, and New Insights. doi:10.1109/tpami.2025.3610085, 2025. \url{https://doi.org/10.1109/tpami.2025.3610085}
\bibitem{ref68} Adversarial Example Detection for DNN Models: A Review and Experimental Comparison. arXiv:2105.00203, 2021. \url{https://arxiv.org/abs/2105.00203}
\bibitem{ref69} Deep Image Restoration Model: A Defense Method Against Adversarial Attacks. doi:10.32604/cmc.2022.020111, 2021. \url{https://doi.org/10.32604/cmc.2022.020111}
\bibitem{ref70} Adv-Makeup: A New Imperceptible and Transferable Attack on Face Recognition. arXiv:2105.03162, 2021. \url{https://arxiv.org/abs/2105.03162}
\bibitem{ref71} A Survey on Adversarial Deep Learning Robustness in Medical Image Analysis. doi:10.3390/electronics10172132, 2021. \url{https://doi.org/10.3390/electronics10172132}
\bibitem{ref72} Proactive Schemes: A Survey of Adversarial Attacks for Social Good. arXiv:2409.16491, 2024. \url{https://arxiv.org/abs/2409.16491}
\bibitem{ref73} Superpixel Attack - Enhancing Black-Box Adversarial Attack with Image-Driven Division Areas. doi:10.1007/978-981-99-8388-9\_12, 2023. \url{https://doi.org/10.1007/978-981-99-8388-9_12}
\bibitem{ref74} Adversarial Examples in the Physical World: A Survey. arXiv:2311.01473, 2023. \url{https://arxiv.org/abs/2311.01473}
\bibitem{ref75} Note on Attacking Object Detectors with Adversarial Stickers. arXiv:1712.08062, 2017. \url{https://arxiv.org/abs/1712.08062}
\bibitem{ref76} Universal Physical Camouflage Attacks on Object Detectors. arXiv:1909.04326, 2019. \url{https://arxiv.org/abs/1909.04326}
\bibitem{ref77} Region-Wise Attack: On Efficient Generation of Robust Physical Adversarial Examples. arXiv:1912.02598, 2019. \url{https://arxiv.org/abs/1912.02598}
\bibitem{ref78} Adversarial Examples: Generation Proposal in the Context of Facial Recognition Systems. arXiv:2404.17760, 2024. \url{https://arxiv.org/abs/2404.17760}
\bibitem{ref79} Robust and Natural Physical Adversarial Examples for Object Detectors. doi:10.1016/j.jisa.2020.102694, 2021. \url{https://doi.org/10.1016/j.jisa.2020.102694}
\bibitem{ref80} Dual Attention Suppression Attack: Generate Adversarial Camouflage in Physical World. doi:10.1109/cvpr46437.2021.00846, 2021. \url{https://doi.org/10.1109/cvpr46437.2021.00846}
\bibitem{ref81} Benchmarking Adversarial Patch Against Aerial Detection. doi:10.1109/tgrs.2022.3225306, 2022. \url{https://doi.org/10.1109/tgrs.2022.3225306}
\bibitem{ref82} Adversarial Examples: Attacks and Defenses for Deep Learning. arXiv:1712.07107, 2017. \url{https://arxiv.org/abs/1712.07107}
\bibitem{ref83} POSES: Patch Optimization Strategies for Efficiency and Stealthiness Using eXplainable AI. doi:10.1109/access.2025.3555044, 2025. \url{https://doi.org/10.1109/access.2025.3555044}
\bibitem{ref84} MVPatch: More Vivid Patch for Adversarial Camouflaged Attacks on Object Detectors in the Physical World. arXiv:2312.17431, 2023. \url{https://arxiv.org/abs/2312.17431}
\bibitem{ref85} PhysPatch: A Physically Realizable and Transferable Adversarial Patch Attack for Multimodal Large Language Models-based Autonomous Driving Systems. arXiv:2508.05167, 2025. \url{https://arxiv.org/abs/2508.05167}
\bibitem{ref86} Enhancing the Transferability of Adversarial Patch via Alternating Minimization. doi:10.1007/s44196-024-00617-2, 2024. \url{https://doi.org/10.1007/s44196-024-00617-2}
\bibitem{ref87} Simultaneously Optimizing Perturbations and Positions for Black-box Adversarial Patch Attacks. arXiv:2212.12995, 2022. \url{https://arxiv.org/abs/2212.12995}
\bibitem{ref88} Generating Adversarial yet Inconspicuous Patches with a Single Image. arXiv:2009.09774, 2020. \url{https://arxiv.org/abs/2009.09774}
\bibitem{ref89} Sparse patches adversarial attacks via extrapolating point-wise information. arXiv:2411.16162, 2024. \url{https://arxiv.org/abs/2411.16162}
\bibitem{ref90} Benchmarking Adversarial Patch Selection and Location. arXiv:2508.01676, 2025. \url{https://arxiv.org/abs/2508.01676}
\bibitem{ref91} Adversarial Sticker: A Stealthy Attack Method in the Physical World. doi:10.1109/tpami.2022.3176760, 2022. \url{https://doi.org/10.1109/tpami.2022.3176760}
\bibitem{ref92} Fast Proxies for LLM Robustness Evaluation. arXiv:2502.10487, 2025. \url{https://arxiv.org/abs/2502.10487}
\bibitem{ref93} PADetBench: Towards Benchmarking Physical Attacks against Object Detection. arXiv:2408.09181, 2024. \url{https://arxiv.org/abs/2408.09181}
\bibitem{ref94} Benchmarking the Physical-world Adversarial Robustness of Vehicle Detection. arXiv:2304.05098, 2023. \url{https://arxiv.org/abs/2304.05098}
\bibitem{ref95} Local Intrinsic Dimensionality Signals Adversarial Perturbations. arXiv:2109.11803, 2021. \url{https://arxiv.org/abs/2109.11803}
\bibitem{ref96} Fast Minimum-norm Adversarial Attacks through Adaptive Norm Constraints. arXiv:2102.12827, 2021. \url{https://arxiv.org/abs/2102.12827}
\bibitem{ref97} LanCe: A Comprehensive and Lightweight CNN Defense Methodology against Physical Adversarial Attacks on Embedded Multimedia Applications. doi:10.1109/asp-dac47756.2020.9045584, 2020. \url{https://doi.org/10.1109/asp-dac47756.2020.9045584}
\bibitem{ref98} Enhancing Security Against Adversarial Attacks Using Robust Machine Learning. doi:10.35940/ijaent.a0485.12010125, 2025. \url{https://doi.org/10.35940/ijaent.a0485.12010125}
\bibitem{ref99} Breaking the Illusion: Real-world Challenges for Adversarial Patches in Object Detection. arXiv:2410.19863, 2024. \url{https://arxiv.org/abs/2410.19863}
\bibitem{ref100} Adversarial robustness assessment: Why in evaluation both L0 and L$\infty$ attacks are necessary. doi:10.1371/journal.pone.0265723, 2022. \url{https://doi.org/10.1371/journal.pone.0265723}
\bibitem{ref101} Tit-for-Tat: Safeguarding Large Vision-Language Models Against Jailbreak Attacks via Adversarial Defense. arXiv:2503.11619, 2025. \url{https://arxiv.org/abs/2503.11619}
\bibitem{ref102} Effectiveness assessment of recent large vision-language models. doi:10.1007/s44267-024-00066-7, 2024. \url{https://doi.org/10.1007/s44267-024-00066-7}
\bibitem{ref103} Immune Defense: A Novel Adversarial Defense Mechanism for Preventing the Generation of Adversarial Examples. arXiv:2303.04502, 2023. \url{https://arxiv.org/abs/2303.04502}
\bibitem{ref104} Safeguarding Vision-Language Models Against Patched Visual Prompt Injectors. arXiv:2405.10529, 2024. \url{https://arxiv.org/abs/2405.10529}
\bibitem{ref105} A New Meta-learning-based Black-box Adversarial Attack: SA-CC. doi:10.1109/ccdc55256.2022.10033574, 2022. \url{https://doi.org/10.1109/ccdc55256.2022.10033574}
\bibitem{ref106} Improved Adversarial Robustness via Logit Regularization Methods. arXiv:1906.03749, 2019. \url{https://arxiv.org/abs/1906.03749}
\bibitem{ref107} Attention, Please! Adversarial Defense via Activation Rectification and Preservation. doi:10.1145/3572843, 2022. \url{https://doi.org/10.1145/3572843}
\bibitem{ref108} Reinforced Embodied Active Defense: Exploiting Adaptive Interaction for Robust Visual Perception in Adversarial 3D Environments. doi:10.1109/tpami.2025.3585726, 2025. \url{https://doi.org/10.1109/tpami.2025.3585726}
\bibitem{ref109} Meta Invariance Defense Towards Generalizable Robustness to Unknown Adversarial Attacks. doi:10.1109/tpami.2024.3385745, 2024. \url{https://doi.org/10.1109/tpami.2024.3385745}
\bibitem{ref110} On Using Certified Training towards Empirical Robustness. arXiv:2410.01617, 2024. \url{https://arxiv.org/abs/2410.01617}
\bibitem{ref111} Recent Advances in Adversarial Training for Adversarial Robustness. arXiv:2102.01356, 2021. \url{https://arxiv.org/abs/2102.01356}
\bibitem{ref112} LAS-AT: Adversarial Training with Learnable Attack Strategy. doi:10.1109/cvpr52688.2022.01304, 2022. \url{https://doi.org/10.1109/cvpr52688.2022.01304}
\bibitem{ref113} Calibrated Adversarial Training. arXiv:2110.00623, 2021. \url{https://arxiv.org/abs/2110.00623}
\bibitem{ref114} Adversarial Distributional Training for Robust Deep Learning. arXiv:2002.05999, 2020. \url{https://arxiv.org/abs/2002.05999}
\bibitem{ref115} Raising the Bar for Certified Adversarial Robustness with Diffusion Models. arXiv:2305.10388, 2023. \url{https://arxiv.org/abs/2305.10388}
\bibitem{ref116} Towards Certified Probabilistic Robustness with High Accuracy. arXiv:2309.00879, 2023. \url{https://arxiv.org/abs/2309.00879}
\bibitem{ref117} Existence and Minimax Theorems for Adversarial Surrogate Risks in Binary Classification. arXiv:2206.09098, 2022. \url{https://arxiv.org/abs/2206.09098}
\bibitem{ref118} Adversarial Surrogate Risk Bounds for Binary Classification. arXiv:2506.09348, 2025. \url{https://arxiv.org/abs/2506.09348}
\bibitem{ref119} Dynamic Label Adversarial Training for Deep Learning Robustness Against Adversarial Attacks. arXiv:2408.13102, 2024. \url{https://arxiv.org/abs/2408.13102}
\bibitem{ref120} How robust accuracy suffers from certified training with convex relaxations. arXiv:2306.06995, 2023. \url{https://arxiv.org/abs/2306.06995}
\bibitem{ref121} Adversarial Machine Learning at Scale. arXiv:1611.01236, 2016. \url{https://arxiv.org/abs/1611.01236}
\bibitem{ref122} Label Smoothing and Logit Squeezing: A Replacement for Adversarial Training?. arXiv:1910.11585, 2019. \url{https://arxiv.org/abs/1910.11585}
\bibitem{ref123} Game-Theoretic Defenses for Robust Conformal Prediction Against Adversarial Attacks in Medical Imaging. arXiv:2411.04376, 2024. \url{https://arxiv.org/abs/2411.04376}
\bibitem{ref124} Introducing Adaptive Continuous Adversarial Training (ACAT) to Enhance ML Robustness. arXiv:2403.10461, 2024. \url{https://arxiv.org/abs/2403.10461}
\bibitem{ref125} VCAT: Vulnerability-aware and Curiosity-driven Adversarial Training for Enhancing Autonomous Vehicle Robustness. arXiv:2409.12997, 2024. \url{https://arxiv.org/abs/2409.12997}
\bibitem{ref126} Adversarial Machine Learning: Attacks, Defenses, and Open Challenges. arXiv:2502.05637, 2025. \url{https://arxiv.org/abs/2502.05637}
\bibitem{ref127} Adversarial example defense based on image reconstruction. doi:10.7717/peerj-cs.811, 2021. \url{https://doi.org/10.7717/peerj-cs.811}
\bibitem{ref128} Defending against Adversarial Attacks using Digital Image Processing. doi:10.1088/1742-6596/2577/1/012016, 2023. \url{https://doi.org/10.1088/1742-6596/2577/1/012016}
\bibitem{ref129} Detecting Adversarial Examples via Key-based Network. arXiv:1806.00580, 2018. \url{https://arxiv.org/abs/1806.00580}
\bibitem{ref130} Feature Losses for Adversarial Robustness. arXiv:1912.04497, 2019. \url{https://arxiv.org/abs/1912.04497}
\bibitem{ref131} Preprocessors Matter! Realistic Decision-Based Attacks on Machine Learning Systems. arXiv:2210.03297, 2022. \url{https://arxiv.org/abs/2210.03297}
\bibitem{ref132} Adversarial Purification through Representation Disentanglement. arXiv:2110.07801, 2021. \url{https://arxiv.org/abs/2110.07801}
\bibitem{ref133} Enhanced Attacks on Defensively Distilled Deep Neural Networks. arXiv:1711.05934, 2017. \url{https://arxiv.org/abs/1711.05934}
\bibitem{ref134} DNNShield: Dynamic Randomized Model Sparsification, A Defense Against Adversarial Machine Learning. arXiv:2208.00498, 2022. \url{https://arxiv.org/abs/2208.00498}
\bibitem{ref135} PCA improves the adversarial robustness of neural networks. doi:10.14428/esann/2022.es2022-96, 2022. \url{https://doi.org/10.14428/esann/2022.es2022-96}
\bibitem{ref136} Deep Latent Defence. arXiv:1910.03916, 2019. \url{https://arxiv.org/abs/1910.03916}
\bibitem{ref137} Robust Adversarial Defense by Tensor Factorization. doi:10.1109/icmla58977.2023.00050, 2023. \url{https://doi.org/10.1109/icmla58977.2023.00050}
\bibitem{ref138} A Self-supervised Approach for Adversarial Robustness. doi:10.1109/cvpr42600.2020.00034, 2020. \url{https://doi.org/10.1109/cvpr42600.2020.00034}
\bibitem{ref139} AdvAttackVis: An Adversarial Attack Visualization System for Deep Neural Networks. doi:10.14569/ijacsa.2024.0150538, 2024. \url{https://doi.org/10.14569/ijacsa.2024.0150538}
\bibitem{ref140} A Framework for Evaluating Vision-Language Model Safety: Building Trust in AI for Public Sector Applications. arXiv:2502.16361, 2025. \url{https://arxiv.org/abs/2502.16361}
\bibitem{ref141} Adversarial Robustness Assessment: Why both \$L\_0\$ and \$L\_\textbackslash\{\}infty\$ Attacks Are Necessary. arXiv:1906.06026, 2019. \url{https://arxiv.org/abs/1906.06026}
\bibitem{ref142} On the Efficacy of Metrics to Describe Adversarial Attacks. arXiv:2301.13028, 2023. \url{https://arxiv.org/abs/2301.13028}
\bibitem{ref143} Stronger and Faster Wasserstein Adversarial Attacks. arXiv:2008.02883, 2020. \url{https://arxiv.org/abs/2008.02883}
\bibitem{ref144} A Framework for Verification of Wasserstein Adversarial Robustness. arXiv:2110.06816, 2021. \url{https://arxiv.org/abs/2110.06816}
\bibitem{ref145} Reliable evaluation of adversarial robustness with an ensemble of diverse parameter-free attacks. arXiv:2003.01690, 2020. \url{https://arxiv.org/abs/2003.01690}
\bibitem{ref146} Orthogonal Deep Models as Defense Against Black-Box Attacks. doi:10.1109/access.2020.3005961, 2020. \url{https://doi.org/10.1109/access.2020.3005961}
\bibitem{ref147} Adversarial training may be a double-edged sword. arXiv:2107.11671, 2021. \url{https://arxiv.org/abs/2107.11671}
\bibitem{ref148} A Survey of Neural Network Robustness Assessment in Image Recognition. arXiv:2404.08285, 2024. \url{https://arxiv.org/abs/2404.08285}
\bibitem{ref149} Defense-friendly Images in Adversarial Attacks: Dataset and Metrics for Perturbation Difficulty. doi:10.1109/wacv48630.2021.00060, 2021. \url{https://doi.org/10.1109/wacv48630.2021.00060}
\bibitem{ref150} The Effects of Image Distribution and Task on Adversarial Robustness. arXiv:2102.10534, 2021. \url{https://arxiv.org/abs/2102.10534}
\bibitem{ref151} An Explainable Adversarial Robustness Metric for Deep Learning Neural Networks. arXiv:1806.01477, 2018. \url{https://arxiv.org/abs/1806.01477}
\bibitem{ref152} BARReL: Bottleneck Attention for Adversarial Robustness in Vision-Based Reinforcement Learning. arXiv:2208.10481, 2022. \url{https://arxiv.org/abs/2208.10481}
\bibitem{ref153} Interpretable Computer Vision Models through Adversarial Training: Unveiling the Robustness-Interpretability Connection. arXiv:2307.02500, 2023. \url{https://arxiv.org/abs/2307.02500}
\bibitem{ref154} Attack-SAM: Towards Attacking Segment Anything Model With Adversarial Examples. arXiv:2305.00866, 2023. \url{https://arxiv.org/abs/2305.00866}
\bibitem{ref155} Evaluating Adversarial Robustness on Document Image Classification. arXiv:2304.12486, 2023. \url{https://arxiv.org/abs/2304.12486}
\bibitem{ref156} A Comprehensive Evaluation Framework for Deep Model Robustness. arXiv:2101.09617, 2021. \url{https://arxiv.org/abs/2101.09617}
\bibitem{ref157} On the Relationship Between Adversarial Robustness and Decision Region in Deep Neural Network. arXiv:2207.03400, 2022. \url{https://arxiv.org/abs/2207.03400}
\bibitem{ref158} Adversarial Robustness in Financial Machine Learning: Defenses, Economic Impact, and Governance Evidence. arXiv:2512.15780, 2025. \url{https://arxiv.org/abs/2512.15780}
\bibitem{ref159} Ti-Patch: Tiled Physical Adversarial Patch for no-reference video quality metrics. arXiv:2404.09961, 2024. \url{https://arxiv.org/abs/2404.09961}
\bibitem{ref160} Quantifying the robustness of deep multispectral segmentation models against natural perturbations and data poisoning. doi:10.1117/12.2663498, 2023. \url{https://doi.org/10.1117/12.2663498}
\bibitem{ref161} Enhancing Cross-task Transferability of Adversarial Examples with Dispersion Reduction. arXiv:1905.03333, 2019. \url{https://arxiv.org/abs/1905.03333}
\bibitem{ref162} Adversarial Attack Vulnerability of Medical Image Analysis Systems: Unexplored Factors. doi:10.1016/j.media.2021.102141, 2021. \url{https://doi.org/10.1016/j.media.2021.102141}
\bibitem{ref163} A Light Recipe to Train Robust Vision Transformers. doi:10.1109/satml54575.2023.00024, 2023. \url{https://doi.org/10.1109/satml54575.2023.00024}
\bibitem{ref164} Robust Models are less Over-Confident. arXiv:2210.05938, 2022. \url{https://arxiv.org/abs/2210.05938}
\bibitem{ref165} PAIF: Perception-Aware Infrared-Visible Image Fusion for Attack-Tolerant Semantic Segmentation. doi:10.1145/3581783.3611928, 2023. \url{https://doi.org/10.1145/3581783.3611928}
\bibitem{ref166} FACESEC: A Fine-grained Robustness Evaluation Framework for Face Recognition Systems. doi:10.1109/cvpr46437.2021.01305, 2021. \url{https://doi.org/10.1109/cvpr46437.2021.01305}
\bibitem{ref167} Holistic Approach to Measure Sample-level Adversarial Vulnerability and its Utility in Building Trustworthy Systems. doi:10.1109/cvprw56347.2022.00479, 2022. \url{https://doi.org/10.1109/cvprw56347.2022.00479}
\bibitem{ref168} How Stealthy is Stealthy? Studying the Efficacy of Black-Box Adversarial Attacks in the Real World. doi:10.1007/978-3-031-92886-4\_10, 2025. \url{https://doi.org/10.1007/978-3-031-92886-4_10}
\bibitem{ref169} Benchmarking the Spatial Robustness of DNNs via Natural and Adversarial Localized Corruptions. doi:10.1016/j.patcog.2025.112412, 2025. \url{https://doi.org/10.1016/j.patcog.2025.112412}
\bibitem{ref170} Augmenting Model Robustness with Transformation-Invariant Attacks. arXiv:1901.11188, 2019. \url{https://arxiv.org/abs/1901.11188}
\bibitem{ref171} ROBY: Evaluating the Robustness of a Deep Model by its Decision Boundaries. arXiv:2012.10282, 2020. \url{https://arxiv.org/abs/2012.10282}
\bibitem{ref172} Provable Adversarial Robustness for Fractional Lp Threat Models. arXiv:2203.08945, 2022. \url{https://arxiv.org/abs/2203.08945}
\bibitem{ref173} Evaluating Model Robustness Using Adaptive Sparse L0 Regularization. arXiv:2408.15702, 2024. \url{https://arxiv.org/abs/2408.15702}
\bibitem{ref174} Perceptual Evaluation of Adversarial Attacks for CNN-based Image Classification. doi:10.1109/qomex.2019.8743213, 2019. \url{https://doi.org/10.1109/qomex.2019.8743213}
\bibitem{ref175} Adversarial Attacks and Defense Mechanisms in Machine Learning: A Structured Review of Methods, Domains, and Open Challenges. doi:10.1109/access.2025.3624409, 2025. \url{https://doi.org/10.1109/access.2025.3624409}
\bibitem{ref176} Benchmarking adversarial attacks and defenses for time-series data. doi:10.1007/978-3-030-63836-8\_45, 2020. \url{https://doi.org/10.1007/978-3-030-63836-8_45}
\bibitem{ref177} Provable robustness against all adversarial \$l\_p\$-perturbations for \$p\textbackslash\{\}geq 1\$. arXiv:1905.11213, 2019. \url{https://arxiv.org/abs/1905.11213}
\bibitem{ref178} Should Adversarial Attacks Use Pixel p-Norm?. arXiv:1906.02439, 2019. \url{https://arxiv.org/abs/1906.02439}
\bibitem{ref179} Imbalanced gradients: a subtle cause of overestimated adversarial robustness. doi:10.1007/s10994-023-06328-7, 2023. \url{https://doi.org/10.1007/s10994-023-06328-7}
\bibitem{ref180} Examining the Human Perceptibility of Black-Box Adversarial Attacks on Face Recognition. arXiv:2107.09126, 2021. \url{https://arxiv.org/abs/2107.09126}
\bibitem{ref181} Lightweight Lipschitz Margin Training for Certified Defense against Adversarial Examples. arXiv:1811.08080, 2018. \url{https://arxiv.org/abs/1811.08080}
\bibitem{ref182} FORTRESS: Frontier Risk Evaluation for National Security and Public Safety. arXiv:2506.14922, 2025. \url{https://arxiv.org/abs/2506.14922}
\bibitem{ref183} Measuring the Robustness of Reference-Free Dialogue Evaluation Systems. arXiv:2501.06728, 2025. \url{https://arxiv.org/abs/2501.06728}
\bibitem{ref184} Perceptually Constrained Adversarial Attacks. arXiv:2102.07140, 2021. \url{https://arxiv.org/abs/2102.07140}
\bibitem{ref185} Lp-norm Distortion-Efficient Adversarial Attack. arXiv:2407.03115, 2024. \url{https://arxiv.org/abs/2407.03115}
\bibitem{ref186} Adversarial Defense Via Local Flatness Regularization. doi:10.1109/icip40778.2020.9191346, 2020. \url{https://doi.org/10.1109/icip40778.2020.9191346}
\bibitem{ref187} The Pitfalls and Promise of Conformal Inference Under Adversarial Attacks. arXiv:2405.08886, 2024. \url{https://arxiv.org/abs/2405.08886}
\bibitem{ref188} Regularization properties of adversarially-trained linear regression. arXiv:2310.10807, 2023. \url{https://arxiv.org/abs/2310.10807}
\bibitem{ref189} Deep-Attack over the Deep Reinforcement Learning. arXiv:2205.00807, 2022. \url{https://arxiv.org/abs/2205.00807}
\bibitem{ref190} ON THE UTILITY OF CONDITIONAL GENERATION BASED MUTUAL INFORMATION FOR CHARACTERIZING ADVERSARIAL SUBSPACES. doi:10.1109/globalsip.2018.8646527, 2018. \url{https://doi.org/10.1109/globalsip.2018.8646527}
\bibitem{ref191} RobustBench: a standardized adversarial robustness benchmark. arXiv:2010.09670, 2020. \url{https://arxiv.org/abs/2010.09670}
\bibitem{ref192} ImageNet-Patch: A Dataset for Benchmarking Machine Learning Robustness against Adversarial Patches. doi:10.1016/j.patcog.2022.109064, 2022. \url{https://doi.org/10.1016/j.patcog.2022.109064}
\bibitem{ref193} NIC-RobustBench: A Comprehensive Open-Source Toolkit for Neural Image Compression and Robustness Analysis. arXiv:2506.19051, 2025. \url{https://arxiv.org/abs/2506.19051}
\bibitem{ref194} MAA: Meticulous Adversarial Attack against Vision-Language Pre-trained Models. arXiv:2502.08079, 2025. \url{https://arxiv.org/abs/2502.08079}
\bibitem{ref195} PROSAC: Provably Safe Certification for Machine Learning Models under Adversarial Attacks. arXiv:2402.02629, 2024. \url{https://arxiv.org/abs/2402.02629}
\bibitem{ref196} Adversarial Attacks on Multi-task Visual Perception for Autonomous Driving. doi:10.2352/j.imagingsci.technol.2021.65.6.060408, 2021. \url{https://doi.org/10.2352/j.imagingsci.technol.2021.65.6.060408}
\bibitem{ref197} Adversarial Attack and Defense of YOLO Detectors in Autonomous Driving Scenarios. arXiv:2202.04781, 2022. \url{https://arxiv.org/abs/2202.04781}
\bibitem{ref198} Dirty Road Can Attack: Security of Deep Learning based Automated Lane Centering under Physical-World Attack. arXiv:2009.06701, 2020. \url{https://arxiv.org/abs/2009.06701}
\bibitem{ref199} Robust Deep Reinforcement Learning for Security and Safety in Autonomous Vehicle Systems. arXiv:1805.00983, 2018. \url{https://arxiv.org/abs/1805.00983}
\bibitem{ref200} Sensor Adversarial Traits: Analyzing Robustness of 3D Object Detection Sensor Fusion Models. arXiv:2109.06363, 2021. \url{https://arxiv.org/abs/2109.06363}
\bibitem{ref201} PhysGAN: Generating Physical-World-Resilient Adversarial Examples for Autonomous Driving. doi:10.1109/cvpr42600.2020.01426, 2020. \url{https://doi.org/10.1109/cvpr42600.2020.01426}
\bibitem{ref202} Fooling Detection Alone is Not Enough: First Adversarial Attack against Multiple Object Tracking. arXiv:1905.11026, 2019. \url{https://arxiv.org/abs/1905.11026}
\bibitem{ref203} D4+: Emergent Adversarial Driving Maneuvers with Approximate Functional Optimization. arXiv:2505.13942, 2025. \url{https://arxiv.org/abs/2505.13942}
\bibitem{ref204} Ad2Attack: Adaptive Adversarial Attack on Real-Time UAV Tracking. arXiv:2203.01516, 2022. \url{https://arxiv.org/abs/2203.01516}
\bibitem{ref205} SADA: Semantic Adversarial Diagnostic Attacks for Autonomous Applications. arXiv:1812.02132, 2018. \url{https://arxiv.org/abs/1812.02132}
\bibitem{ref206} A Taxonomy of System-Level Attacks on Deep Learning Models in Autonomous Vehicles. doi:10.1145/3769009, 2025. \url{https://doi.org/10.1145/3769009}
\bibitem{ref207} Physically Realizable Adversarial Examples for LiDAR Object Detection. doi:10.1109/cvpr42600.2020.01373, 2020. \url{https://doi.org/10.1109/cvpr42600.2020.01373}
\bibitem{ref208} Defending Against Person Hiding Adversarial Patch Attack with a Universal White Frame. arXiv:2204.13004, 2022. \url{https://arxiv.org/abs/2204.13004}
\bibitem{ref209} Enhancing Security in Real-Time Video Surveillance: A Deep Learning-Based Remedial Approach for Adversarial Attack Mitigation. doi:10.1109/access.2024.3418614, 2024. \url{https://doi.org/10.1109/access.2024.3418614}
\bibitem{ref210} Real-Time Robust Video Object Detection System Against Physical-World Adversarial Attacks. doi:10.1109/tcad.2023.3305932, 2023. \url{https://doi.org/10.1109/tcad.2023.3305932}
\bibitem{ref211} Adversarial Machine Learning Attacks Against Video Anomaly Detection Systems. doi:10.1109/cvprw56347.2022.00034, 2022. \url{https://doi.org/10.1109/cvprw56347.2022.00034}
\bibitem{ref212} SSMI: How to Make Objects of Interest Disappear without Accessing Object Detectors?. arXiv:2206.10809, 2022. \url{https://arxiv.org/abs/2206.10809}
\bibitem{ref213} AiWatch: A Distributed Video Surveillance System Using Artificial Intelligence and Digital Twins Technologies. doi:10.3390/technologies13050195, 2025. \url{https://doi.org/10.3390/technologies13050195}
\bibitem{ref214} AdvFaces: Adversarial Face Synthesis. doi:10.1109/ijcb48548.2020.9304898, 2020. \url{https://doi.org/10.1109/ijcb48548.2020.9304898}
\bibitem{ref215} AdvFaceGAN: a face dual-identity impersonation attack method based on generative adversarial networks. doi:10.7717/peerj-cs.2904, 2025. \url{https://doi.org/10.7717/peerj-cs.2904}
\bibitem{ref216} Real-World Adversarial Examples involving Makeup Application. arXiv:2109.03329, 2021. \url{https://arxiv.org/abs/2109.03329}
\bibitem{ref217} Detection of Face Recognition Adversarial Attacks. arXiv:1912.02918, 2019. \url{https://arxiv.org/abs/1912.02918}
\bibitem{ref218} A Comprehensive Risk Analysis Method for Adversarial Attacks on Biometric Authentication Systems. doi:10.1109/access.2024.3439741, 2024. \url{https://doi.org/10.1109/access.2024.3439741}
\bibitem{ref219} Detecting Adversarial Faces Using Only Real Face Self-Perturbations. arXiv:2304.11359, 2023. \url{https://arxiv.org/abs/2304.11359}
\bibitem{ref220} ApaNet: adversarial perturbations alleviation network for face verification. doi:10.1007/s11042-022-13641-1, 2022. \url{https://doi.org/10.1007/s11042-022-13641-1}
\bibitem{ref221} Edge Detection with a Preprocessing Approach. doi:10.4236/jsip.2014.54015, 2014. \url{https://doi.org/10.4236/jsip.2014.54015}
\bibitem{ref222} NeRFail: Neural Radiance Fields-Based Multiview Adversarial Attack. doi:10.1609/aaai.v38i19.30113, 2024. \url{https://doi.org/10.1609/aaai.v38i19.30113}
\bibitem{ref223} Blocks. doi:10.1145/3351241, 2019. \url{https://doi.org/10.1145/3351241}
\bibitem{ref224} Multiple Perturbation Attack: Attack Pixelwise Under Different \$\textbackslash\{\}ell\_p\$-norms For Better Adversarial Performance. arXiv:2212.03069, 2022. \url{https://arxiv.org/abs/2212.03069}
\bibitem{ref225} Defending Adversarial Attacks by Correcting logits. arXiv:1906.10973, 2019. \url{https://arxiv.org/abs/1906.10973}
\bibitem{ref226} Adversarial Vision Challenge. doi:10.1007/978-3-030-29135-8\_5, 2019. \url{https://doi.org/10.1007/978-3-030-29135-8_5}
\bibitem{ref227} Towards Adversarial Attack on Vision-Language Pre-training Models. arXiv:2206.09391, 2022. \url{https://arxiv.org/abs/2206.09391}
\bibitem{ref228} Improving Adversarial Robustness via Attention and Adversarial Logit Pairing. doi:10.3389/frai.2021.752831, 2022. \url{https://doi.org/10.3389/frai.2021.752831}
\bibitem{ref229} Multimodal Robust Prompt Distillation for 3D Point Cloud Models. arXiv:2511.21574, 2025. \url{https://arxiv.org/abs/2511.21574}
\bibitem{ref230} CoDefend: Cross-Modal Collaborative Defense via Diffusion Purification and Prompt Optimization. arXiv:2510.11096, 2025. \url{https://arxiv.org/abs/2510.11096}
\bibitem{ref231} Using Multiple Self-Supervised Tasks Improves Model Robustness. arXiv:2204.03714, 2022. \url{https://arxiv.org/abs/2204.03714}
\bibitem{ref232} SEC4SR: A Security Analysis Platform for Speaker Recognition. arXiv:2109.01766, 2021. \url{https://arxiv.org/abs/2109.01766}
\bibitem{ref233} SuperPure: Efficient Purification of Localized and Distributed Adversarial Patches via Super-Resolution GAN Models. arXiv:2505.16318, 2025. \url{https://arxiv.org/abs/2505.16318}
\bibitem{ref234} PATCHOUT: Adversarial Patch Detection and Localization using Semantic Consistency. doi:10.1007/s11063-025-11775-5, 2025. \url{https://doi.org/10.1007/s11063-025-11775-5}
\bibitem{ref235} On the Adversarial Robustness of Camera-based 3D Object Detection. arXiv:2301.10766, 2023. \url{https://arxiv.org/abs/2301.10766}
\end{thebibliography}

\end{document}
